% Does anybody know how was the header before? I just jumped in and I saw the \documentclass being broken. Whoever accidentally broke it, could please fix it to bring it to how it was before? I am fixing it just to make it work, for the time being. Thank you 
\documentclass[twocolumn,preprintnumbers,
  superscriptaddress,
  nofootinbib,
  amsmath,
  amssymb,
  aps,
  colorlinks=true,
  linkcolor=blue,
  urlcolor=blue,
  citecolor=blue,
  anchorcolor=blue
]{revtex4-2}

\usepackage{wdmmanuscript}
\usepackage{hyperref}
\usepackage{fontawesome5}
\usepackage{ulem}
\usepackage{mathtools}

\newcommand{\codelink}[1]{\,\href{#1}{\scalebox{0.65}{\faGithub}}}

\newcommand{\TODO}[1]{\textcolor{red}{\textbf{[TODO: #1]}}}

\newcommand{\ob}[1]{\textcolor{blue}{\textbf{[Ollie: #1]}}}
\newcommand{\gm}[1]{\textcolor{red}{\textbf{[Giorgio: #1]}}}
\newcommand{\avi}[1]{\textcolor{red}{\textbf{[Avi: #1]}}}
\newcommand{\ls}[1]{\textcolor{magenta}{\textbf{[Lore: #1]}}}
\newcommand{\qb}[1]{\textcolor{magenta}{\textbf{[Quentin: #1]}}}

\begin{document}

\preprint{APS/123-QED}

\title{JAX-based codes for WDM-transform \\for use within gravitational wave astronomy}

\author{Quentin Baghi}
\affiliation{Universit\'e Paris Cit\'e, CNRS, Astroparticule et Cosmologie, F-75013 Paris, France}
\author{Ollie Burke}
\affiliation{School of Physics and Astronomy, University of Glasgow,
University Avenue, Glasgow, G12 8QQ, UK}
\affiliation{Laboratoire des 2 Infinis - Toulouse (L2IT-IN2P3),
Université de Toulouse, CNRS, UPS, F-31062 Toulouse Cedex 9, France}
\author{Giorgio Mentasti}
\affiliation{Universit\'e Paris Cit\'e, CNRS, Astroparticule et Cosmologie, F-75013 Paris, France}
\author{Lorenzo Speri}
\affiliation{%
European Space Agency (ESA), European Space Research and Technology Centre (ESTEC), Keplerlaan 1, 2201 AZ Noordwijk, the Netherlands
}
\affiliation{Leiden Observatory, Leiden University, P.O. Box 9513, 2300 RA Leiden, the Netherlands}
\author{Avi Vajpeyi}
\affiliation{Department of Statistics, University of Auckland, 38
Princes St, Auckland, New Zealand}

\begin{abstract}
  The WDM transform has been widely used in gravitational-wave astronomy, yet a
  self-contained, mathematically explicit reference for practitioners remains lacking.
  This is especially true for those wishing to adopt the transform in modern Python
  and JAX inference workflows.
  We present \package, an open-source Python package implementing the
  Wilson--Daubechies--Meyer (WDM) wavelet-packet time-frequency transform for
  gravitational-wave data analysis.
  We document the forward and inverse transform step by step, making explicit the
  packed $(N_t, N_f{+}1)$ coefficient layout, the DC and Nyquist edge-channel
  conventions, the alternating phase factor $C_{nm}$, and the normalization
  conditions required for exact round-trip reconstruction.
  The package supplies NumPy and JAX backends, both round-trip validated to
  floating-point precision, with the JAX backend enabling GPU-accelerated transforms
  of million-point data streams in tens of milliseconds. As a worked example we
  verify that the WDM-domain Whittle likelihood reproduces frequency-domain
  posteriors for a resolved LISA galactic binary under a shared stationary noise
  model, confirming numerical equivalence of the two representations in that
  controlled setting.
\end{abstract}

\maketitle



\section{Introduction}





Time-frequency analysis is widely used in gravitational-wave
astronomy. It enables template-free characterisation of source
signals~\cite{krolak_1996,Virtuoso2024}, instrumental
transients~\cite{robinet2020}, and the joint analysis of
astrophysical signals and terrestrial
disturbances~\cite{Cornish_2015}. Time-frequency
representations also provide compact waveform descriptions and
facilitate the modelling of non-stationary detector noise. Among the
available techniques, discrete wavelet-packet transforms are
particularly well suited to gravitational-wave data analysis: they
furnish localized, orthonormal basis functions capable of representing
transient and chirping signals, and for locally stationary noise they
yield an approximately diagonal noise covariance, greatly simplifying
likelihood evaluation and spectral estimation.

The Wilson--Daubechies--Meyer (WDM) wavelet-packet basis already
has a substantial history in gravitational-wave astronomy.  The fast
Wilson--Daubechies transform was developed for transient
gravitational-wave analysis by Necula, Klimenko, and
Mitselmakher~\cite{Necula_2012}, and WDM representations have been
used in coherent WaveBurst and related burst-search contexts~\cite{2023PhRvD.107f2002S, 2021SoftX..1400678D}.
These earlier papers and implementations established the theoretical
and algorithmic basis of the transform for excess-power and
time-frequency analyses, including its orthonormal structure,
controlled spectral leakage, flexible frequency sub-bands, and
time-delay filters relevant for sky localization.  Subsequent WDM
software and applications by Digman and Cornish further developed fast
forward and inverse transforms, WDM likelihoods, and time-frequency
analyses for LISA sources and backgrounds~\cite{2023ascl.soft07037D,
digman_cornish_bbh_wdm,digman_cornish_tv_galactic_wdm}.  Recent work
by Pearson and Cornish has also shown how WDM-domain methods can be
combined with Bayesian data augmentation to handle data gaps in
long-duration gravitational-wave observations~\cite{pearson_cornish_wdm_gaps}.

Cornish developed a complementary WDM-domain programme for
likelihood-based gravitational-wave inference and waveform generation.
In this formulation, regularly sampled data are mapped to an
orthogonal grid of localized time-frequency pixels; slowly evolving
signals may be represented compactly as tracks on this grid, and
binary waveforms can be generated directly in the WDM domain
\cite{Cornish:2020odn}.  A central statistical motivation is that
locally stationary noise, whose spectrum varies slowly across an
individual time-frequency pixel, has an approximately diagonal
covariance in the WDM domain.  This approximate diagonalization,
analyzed in detail in
Ref.~\cite{cornish_nonstationary_wdm}, provides the WDM-domain
analogue of the familiar diagonal Fourier-domain covariance for
stationary noise and underlies WDM Whittle-type likelihoods. A clear one-to-one correspondence between equations and their implementation is often missing in existing literature and codebases \cite{Cornish_WDM_Transform,Digman_WDMWaveletTransforms,Moore_WDM_GW_wavelets,cWB_software}:
We address this gap here.

The present manuscript does not claim to introduce the WDM transform,
to replace these foundational analyses, or to outperform existing
pipeline-specific implementations.  Its contribution is instead to
make the sampled WDM construction easy to adopt and reproduce in
modern Python-based gravitational-wave workflows.  We give a
self-contained derivation of the forward and inverse transforms,
state the packed $(N_t,N_f+1)$ coefficient layout explicitly, spell
out the DC and Nyquist channel treatment, and fix the normalization,
phase, and indexing conventions used by the implementation.  These
choices are then validated through exact round-trip reconstruction
tests and empirical checks of coefficient decorrelation.  The
accompanying \package\ package provides CPU and GPU-capable NumPy/JAX
backends, so that the same mathematically explicit transform can be
used in conventional analyses and in differentiable inference
workflows.

The motivation for the WDM can be understood by contrasting the two
standard representations of a discrete signal. 
In the time domain one
can identify \emph{when} a feature occurs but not which frequency
band is active, whereas in the Fourier domain one can identify
\emph{which} frequencies are present but not when they appear. Many
gravitational-wave signals and detector artefacts are neither purely
stationary nor purely impulsive, so neither representation alone is
fully satisfactory. 
A short-time Fourier often makes use of compactly-supported sliding windows such as square or Tukey windows. In contrast, WDM employs frequency-localized atoms whose corresponding time-domain representation exhibits algebraically decaying tails proportional to $1/t^3$.

The result is a
uniform time-frequency representation well suited to describing
transients, chirps, gaps, and other localized structure, while still
admitting an exact inverse transform back to the original sampled data.

The remainder of this paper is organized as follows.
Section~\ref{sec:conventions} introduces notation and conventions.
Section~\ref{sec:wdm} develops the WDM transform, including the Meyer
window, the full three-case basis definition covering DC, interior,
and Nyquist channels, the forward and inverse transforms, and the
statistical properties of the coefficients under stationary and
non-stationary noise. Section~\ref{sec:code_package} describes the
\package\ package, presents runtime and accuracy benchmarks, and
demonstrates the transform on astrophysical signals. We conclude in
Section~\ref{sec:conclusions}.

\subsection{Conventions}\label{sec:conventions}
Consider a data stream consisting of $N$ uniformly sampled values of a
\textbf{real} quantity $x$ at times $\{t_n\}_{n=0,\dots,N-1}$ over a
total observation time $T$:
\begin{align}\label{st_definition}
  x[n]=x(t_n)\qquad,\qquad n=0,\dots,N-1\;,
\end{align}
with sampling interval $\Delta t = t_{k+1} - t_{k}$ and $T = N\Delta t$.
All available information about the signal $x[n]$ is encoded in these $N$ samples.
We define the sampling rate and the Nyquist frequency as
\begin{align}
  F_s&\equiv\frac{1}{\Delta t}\equiv\frac{N}{T}\;,\nonumber\\
  f_{\rm max}&=\frac 1 2F_s\;,\nonumber\\
  \Delta f &= 1/T
\end{align}
In that regard, if we set $t_0=0$, we can write $t_n=n\Delta t$.
In the frequency domain, the data stream $\{x_n\}_{n=0}^{N-1}$ can be
discrete-Fourier-transformed (DFT) to a set of other $N$ complex numbers as
\begin{equation}
  \label{eq:dft-definition}
  \tilde x[l] = \sum_{k = 0}^{N-1} x[k]e^{-2\pi i l k/N}
\end{equation}
for Fourier frequencies $f_{l} = \{-N\Delta f/2, \ldots, -\Delta f,
0, \Delta f, \ldots, (N/2 - 1)\Delta f\}$, i.e.\ the zero-frequency
bin is centred in the array. Bold faced quantities $\boldsymbol{A}$ stand for vector/matrix quantities with components $A_{ij}$. Dagger based operations $\boldsymbol{A}^{\dagger}$ stand for the conjugate transpose of the matrix $\boldsymbol{A}$.

We assume that the output of a gravitational wave detector, $x(t)$,
is expressed as a linear combination of a signal, $h(t; \boldsymbol{\theta})$,
determined by a finite set of (unknown) parameters, $\boldsymbol{\theta}$, and
instrumental noise, $n(t)$. If we ignore the presence of calibration
errors~\cite{Savalle_2022}, the content of a single data stream,
i.e., one output channel from one detector, can be written in the
time and frequency domain as:
\begin{align}
  {x}[k] & = h[k;\boldsymbol{\theta}] + n[k],\\
  \tilde x[l]& = \tilde{h}[l;\boldsymbol{\theta}]+ \tilde{n}[l]
\end{align}
where the tilde indicates the Fourier transform.

\section{The Wilson--Daubechies--Meyer (WDM) Transform}\label{sec:wdm}

\subsection{Terminology}
The name Wilson--Daubechies--Meyer combines several pieces of
time-frequency analysis terminology. A standard discrete wavelet
transform usually refers to a multiresolution, dyadic decomposition
(``dyadic'' meaning organized by factors of two, so that each
successive scale changes the time and frequency resolution by powers
of two). Such transforms represent low-frequency structure with broad
time support and fine frequency resolution, and high-frequency
structure with narrower time support and coarser frequency
resolution. A \emph{wavelet packet} (or \emph{wavelet-wavepacket})
transform is more general: it allows a broader class of admissible
tilings of the time-frequency plane. The object studied in this paper
is the Wilson--Daubechies--Meyer (WDM) wavelet-packet transform; for
brevity, we refer to it as the WDM transform throughout. Unlike a
standard dyadic wavelet transform, the WDM transform used here has a
uniform rectangular tiling, so that all pixels have the same time
width $\Delta T$ and frequency width $\Delta F$, with
$\Delta T\,\Delta F = 1/2$ under the one-sided frequency convention
used below~\cite{Daubechies:1991wv}.%

The naming of the transform reflects two distinct ingredients: the
\emph{Wilson}-basis pairing of positive- and negative-frequency atoms,
introduced together with the phase factor, and the smooth
\emph{Daubechies--Meyer} frequency window used to
construct those atoms, all of which introduced together with its closed form in
Sec.~\ref{sec:wdm} below. 

The WDM transform maps the original $N$ samples to a packed
coefficient grid with $N_t$ time bins and $N_f+1$ stored frequency
channels. The transform size is still set by
\begin{align}\label{NtNf_def}
  N=N_t\,N_f\;,
\end{align}
and the stored channels are indexed by $m=0,\dots,N_f$, with $m=0$
the DC edge channel, $m=1,\dots,N_f-1$ the interior channels, and
$m=N_f$ the Nyquist edge channel. The underlying idea is to partition the data stream into $N_t$ time
segments and to perform a frequency domain analysis of each one. Here $N$ is
fixed by the data, while $N_t$ is a free user choice subject only to
$1 \le N_t \le N$ (and consequently $1 \le N_f \le N$).
Given that choice we can define the spacing in both time and frequency with respect to the underlying tiling with area $N = N_{t}N_{f}$
\begin{align}\label{DTDF}
  \Delta T&\equiv\frac{T_{\text{obs}}}{N_t}\; \qquad \text{grid time
  spacing},\nonumber\\
  \Delta F&\equiv\frac{f_{\rm max}}{N_f} \qquad \text{grid frequency spacing}\,.
\end{align}
One immediately sees that,
\begin{align}
  \Delta T &= \frac{T_{\text{obs}}}{N_{t}} = \frac{N\Delta t}{N_{t}}
  = N_{f}\Delta t \\
  \Delta F &= \frac{f_{\text{max}}}{N_{f}} = \frac{1}{2N_{f}\Delta t}=\frac{N_t\Delta f}{2} = \frac{1}{2\Delta T}
\end{align}
This final expression shows how the Nyquist theorem still holds. Once
you fix $N_t$ (or equivalently $N_f$, $\Delta T$ or $\Delta F$) the
grid is uniquely identified. Figure~\ref{fig:wdm_transform}
summarizes the packed layout, including the DC and Nyquist edge channels and the alternating phase factor used for the interior coefficients. Note that $N_t$ and $N_f$ are required to be even — they need not be powers of two.
Relaxing the evenness requirement is left to future work.


\begin{figure*}[ht]
    \centering
    \includegraphics[width=1.0\linewidth]{figures/wdm_overview.pdf}%
    \caption{
        Schematic illustration of the WDM forward transform.
        (a)~The input DFT spectrum $\tilde{x}[\ell]$ with channel $m$
        spanning $N_t$ bins centred at $mN_t/2$.
        (b)~The Meyer window $\tilde{\varphi}_m[\ell]$ shifted across
        $N_f$ channels.
        (c)~The time-domain atom $x_m[n]$ obtained by applying the
        inverse DFT to the windowed channel.
        (d)~The packed coefficient array $W_{nm}$, with each cell
        $w_{nm} = \mathrm{Re}[C^*_{nm}\,x_m[n]]$; the phase factor
        $C_{nm}$ alternates between $1$ and $i$ on a checkerboard
        pattern (see legend).
    }
    \label{fig:wdm_transform}
    \script{wdm_overview_panels.py}
\end{figure*}

\subsection{Forward transform}
For a given DFT frequency series $\tilde x[l]$, the WDM transform represents the data through the WDM coefficients $w_{nm}$
\begin{align}\label{wnm_def}
w_{nm}&=\sum_{l=-N/2}^{N/2-1}\tilde x[l]\tilde g^*_{nm}[l]\,,\codelink{https://github.com/pywavelet/wdm_transform/blob/v0.5.0/src/wdm_transform/transforms/xp_backend.py\#L225}
\end{align}
where $n$ corresponds to the time index and $m$ corresponds to the
frequency-channel.

The Wilson-basis elements $\tilde{g}$ with Daubechies-Meyer frequency-domain window $\tilde{\varphi}$ are defined in closed form in the discrete frequency domain below. For DC ($m=0$), interior channels ($0 < m < N_f$), and
Nyquist ($m=N_f$) we have that:
\begin{widetext}
\begin{align}
\label{eq:gnm_full_discrete}
\text{\textbf{Wilson Basis:}}\quad \tilde{g}_{nm}[l] &=
\dfrac{1}{\sqrt{2}}\begin{cases}
e^{-4\pi i n l/N_t}\, \tilde{\varphi}[l], & m=0, \\[6pt]
e^{-2\pi i n l/N_t}\Big(C_{nm}\,\tilde{\varphi}[l - mN_t/2] + C_{nm}^*\,\tilde{\varphi}[l + mN_t/2]\Big), & 0 < m < N_f, \\[6pt]
e^{-4\pi i n l/N_t}\Big(\tilde{\varphi}[l - N/2] + \tilde{\varphi}[l + N/2]\Big), & m = N_f, 
\end{cases}
\codelink{https://github.com/pywavelet/wdm_transform/blob/v0.5.0/src/wdm_transform/windows.py\#L129}\\
\text{\textbf{Daubechies--Meyer:}} \quad  \tilde\varphi[l]&=\sqrt{\frac{2}{N_t}}\begin{cases}
1 & |l_r|<A \\
\cos \left[\frac{\pi}{2}\,\left(
\frac{|l_r|-A}{B}\right)\right] & A \leq|l_r| < A+B\,, \\
0 & \text{otherwise}\,,
\end{cases} \qquad l_r =l\frac{2}{N_t}\,,\codelink{https://github.com/pywavelet/wdm_transform/blob/v0.5.0/src/wdm_transform/windows.py\#L71} \label{eq:DM_window_discrete}\\
  C_{nm}&=e^{\frac{i\pi}{4}(1-(-1)^{n+m})} =
  \begin{cases} 1 & n+m \text{ even,} \\ i & n+m \text{ odd.}\label{eq:c_nm_def}
\end{cases}\codelink{https://github.com/pywavelet/wdm_transform/blob/v0.5.0/src/wdm_transform/windows.py\#L58}
\end{align}
\begin{equation}\label{eq:A_B_condition}
B = 1 - 2A \,, \quad \text{for}\ A\in(0,1/2)\,. 
\end{equation}
\end{widetext}
Equation~\eqref{eq:A_B_condition} gives freedom to the user to define the properties of the Daubechies--Meyer window defined in Eq.~\eqref{eq:DM_window_discrete}.

From this point onwards, we will provide an interpretation for each of the equations \eqref{eq:gnm_full_discrete} -- \eqref{eq:A_B_condition}. We begin by discussing the DC, Interior, Nyquist split of the basis functions $\tilde{g}$, an important feature of our codebase, and the practical and theoretical consequences of the Daubechies--Meyer window $\tilde{\varphi}$.

% This orthogonality is the key reason that the WDM formalism is adept for treating stationary noise processes for reasons that will be discussed later.

From Eq.\eqref{eq:c_nm_def}, the interior channels carry an alternating phase factor $C_{nm}$, while the DC and Nyquist edge channels have a doubled time-shift exponent and no $C_{nm}$ modulation. The phase factor $C_{nm}$ is the
``\emph{Wilson}'' ingredient of the WDM construction: it combines the
positive- and negative-frequency windowed Fourier atoms centered at $mN_t/2$ into a single real, orthogonal pair, which is what
makes the interior basis real-valued and Wilson-orthogonal, to be discussed later in this section, in the
discrete setting. 

These edge channels $m = 0$ and $m = N_f$ in \eqref{eq:gnm_full_discrete} each contribute $N_t$ real degrees of freedom to
the packed grid: the interior block stores $N_t(N_f{-}1)$ real
coefficients via the $C_{nm}$-modulated pairing of the
$\pm mN_t/2$ atoms (the Wilson-basis construction), while the DC
and Nyquist channels each contribute $N_t$ unpaired real coefficients.
The total $N_t(N_f{+}1)$ stored entries therefore represent the same
$N=N_tN_f$ time-domain degrees of freedom, with the edge channels
being exactly $N_t$ real numbers~\cite{Necula_2012,Cornish:2020odn}. 
We refer the reader to Fig.~\ref{fig:wdm_transform} for a schematic of the WDM transform and the packed coefficient layout.  % 

The window $\tilde\varphi$ is the ``\emph{Daubechies--Meyer}''
ingredient of the WDM construction: a smooth, compactly supported
frequency window of the Meyer family, whose taper region overlaps
between adjacent channels in a way that yields exact partition-of-unity
reconstruction. The Meyer family is parametrized by an integer order
$d \geq 1$ that controls the smoothness of the taper: the cosine
argument $(|l_r|-A)/B$ in Eq.~\eqref{eq:DM_window_discrete} is replaced
by the normalized incomplete beta function
\begin{equation}\label{eq:nu_d_def}
  \nu_d(y) = \frac{\int_0^y t^{d-1}(1-t)^{d-1}\,\mathrm{d}t}
                  {\int_0^1 t^{d-1}(1-t)^{d-1}\,\mathrm{d}t}\,,
  \qquad y = \frac{|l_r|-A}{B}\,,
\end{equation}
which satisfies $\nu_d(y) + \nu_d(1-y) = 1$ and steepens the roll-off
as $d$ increases. Setting $d=1$ gives $\nu_1(y)=y$ and recovers the
cosine taper of Eq.~\eqref{eq:DM_window_discrete}. In the present
implementation, and in this work, we use this
$d=1$ cosine-tapered member of the family. We do stress that the parameter $d$ is retained in the Application Programming Interface (API) only for
possible future
generalizations.\footnote{Cornish~\cite{Cornish:2020odn} uses $d=4$,
  which yields a flatter passband and steeper roll-off at the expense
  of a wider time-domain support (the window function must be truncated
  at $\pm q\Delta T$ with $q=16$). The $d=1$ cosine taper used here has
  a gentler roll-off but is more compact in time, requiring a smaller
  $q$. In practice, the choice of $d$ trades spectral leakage
  suppression against temporal compactness of the window; for the
  applications considered in this paper the $d=1$ window provides
  adequate frequency separation while keeping the implementation simple
and the time-domain filter short.} By construction, $\tilde\varphi$ has a limited frequency bandwidth
which is controlled by the parameter $A\in(0,1/2)$ in Eq.\eqref{eq:A_B_condition}. 
The frequency localization of the WDM atoms is illustrated in Fig.~\ref{fig:wdm_transform}(b), where the same two highlighted cells are shown in frequency space.
At each wavelet frequency bin $l$, one sees that the window overlaps with adjacent frequency bins.  

For our implementation we use the default $A =1/3$ and
hence $B = 1/3$, so the window is non-zero for normalised frequencies
$\leq 2/3$. This is a default, not a hard-coded
constraint: $A$ is exposed in \texttt{wdm\_transform} as the
parameter \texttt{a} of the \texttt{WDM} class (with $B=1-2A$ derived
automatically), and any choice with $A\in(0,1/2)$ is admissible.%

A final and essential property of the WDM formalism is the principle of orthogonality. The Wilson basis $\tilde{g}$ defined by Eq.~\eqref{eq:gnm_full_discrete} combined with the Daubechies--Meyer window Eq.~\eqref{eq:DM_window_discrete} with associated constants~\eqref{eq:c_nm_def} provide an orthogonal basis\footnote{Our definitions of the Wilson basis functions \eqref{eq:gnm_full_discrete} carefully include the $m=0$ and $m = N_f$ in order for orthogonality to hold.}, with both of the following equalities proved in Appendix~\ref{app:normality_condition} using Eq.~\eqref{eq:DM_window_discrete} and Eq.~\eqref{eq:c_nm_def}:%
\begin{widetext}
\begin{align}\label{eq:orthonormality}
\sum_{l=-N/2}^{N/2-1}\tilde{g}^*_{nm}[l]\tilde{g}_{pq}[l]&=\delta_{m,q}\begin{cases}
    \frac{1}{2}\left(\delta_{n,p}+\delta_{n,p\pm \frac{N_t}{2}}\right) \qquad m\in\{0,N_f\}\\
    \delta_{n,p} \qquad 0<m<N_f
\end{cases}\nonumber\\
\sum_{n=0}^{N_t-1}\sum_{m=0}^{N_f}\tilde{g}^*_{nm}[l]\tilde{g}_{nm}[l']&=\delta_{ll'}\;,
\end{align}
\end{widetext}

The basis functions here are represented in the frequency domain and it is easily proved they are orthogonal in the time domain.
Orthogonality is one of the key reasons that the WDM formalism is adept for gravitational-wave data analysis, which will be discussed at a later stage of this work.



These are the main ingredients behind the packed WDM grid. A single window shape is shifted across
frequency to define the channel index $m$, and each shifted window is then modulated in time to define the atom indexed by $n$. This is why one WDM coefficient can be read as the strength of one localized time-frequency pattern, rather than as the amplitude of a globally supported sinusoid. 

More concretely, a single coefficient $w_{nm}$ answers the question:
``how much does the data resemble the atom centered at time bin $n$
and frequency bin $m$?'' A large magnitude means that the data
contains a feature with roughly the same time localization and central
frequency. 
Reading the WDM plane coefficient-by-coefficient in this way is often much more
intuitive than interpreting a global Fourier series, because the
index pair $(n,m)$ already carries an approximate time-frequency location.

In the next subsections, we will provide formulae for the $w_{nm}$ coefficients for both the interior, DC and Nyquist channels separately. In our final subsection, we investigate the orthogonal nature of the Wilson-basis.  

\subsubsection{Interior channels \texorpdfstring{$(0 < m < N_f)$}{(0 < m < N_f)}}

Starting from Eq.~\eqref{wnm_def} with the interior-channel basis
element from Eq.~\eqref{eq:gnm_full_discrete}, the forward transform can be
evaluated in the Fourier domain using
\begin{align}
  w_{nm} = \sum_{l=-N/2}^{N/2-1} \tilde{x}[l]\,\tilde{g}_{nm}^*[l]\,,
\end{align}
For a real-valued input signal, $\tilde{x}[-l] = \tilde{x}^*[l]$, and the real-valued window satisfies $\tilde{\varphi}[-l] = \tilde{\varphi}[l]$.
Splitting the sum into positive and negative frequency halves and
using these symmetries, the negative-frequency contribution is the
complex conjugate of the positive-frequency contribution (see
Appendix~\ref{app:forward_derivation} for details), giving
%
\begin{equation}\label{eq:wnm_interior}
\begin{split}
  w_{nm} = \sqrt{2}&\,(-1)^{nm}\,\Re\Bigl[C_{nm}^*\\
  &\times\sum_{l'=-N_t/2}^{N_t/2-1}\tilde{x}[l'{+}mN_t/2]\,e^{2\pi i l'n/N_t}\,\tilde{\varphi}[l']\Bigr],\codelink{https://github.com/pywavelet/wdm_transform/blob/v0.5.0/src/wdm_transform/transforms/xp_backend.py\#L92-L102}
\end{split}
\end{equation}
where $l' = l - mN_t/2$ and $\Delta f = 1/T$.
The expression above can be written compactly as
%
\begin{align}\label{eq:wnm_compact}
  w_{nm} = \sqrt{2}\,(-1)^{nm}\,\Re\,C_{nm}^*\,\tilde x_m[n]\,,
\end{align}
%
where
%
\begin{align}\label{eq:xm_def}
  \tilde x_m[n] = \sum_{l=-N_t/2}^{N_t/2-1} e^{2\pi i l n/N_t}\,\tilde x[l +
  mN_t/2]\,\tilde{\varphi}[l],
\end{align}
%
The quantity
$\tilde x_m[n]$ can be computed using a length-$N_t$ FFT for each channel
$m$, so the total cost of the interior channels is $N_f \times
\mathcal{O}(N_t\log N_t)$.

\subsubsection{DC channel \texorpdfstring{$(m = 0)$}{(m = 0)}}

The DC edge channel has a doubled time-shift exponent and no $C_{nm}$
factor in the basis definition~\eqref{eq:gnm_full_discrete}. Following the
same Fourier-domain evaluation yields
\begin{align}\label{eq:wn0}
  w_{n0} = \sqrt{2\,}\Re\left[\sum_{l=1}^{N_t/2-1}\tilde{x}[l]\,e^{4\pi i l
  n/N_t}\,\tilde{\varphi}[l]\right] + \frac{1}{\sqrt{2}}\tilde{x}[0]\,\tilde{\varphi}[0].\codelink{https://github.com/pywavelet/wdm_transform/blob/v0.5.0/src/wdm_transform/transforms/xp_backend.py\#L80-L90}
\end{align}
Note the factor of $4\pi$ in the exponent (versus $2\pi$ for interior
channels): the DC basis element oscillates at twice the rate in the
time index because it sees both positive and negative frequency
copies of the window centered at $f=0$.

\subsubsection{Nyquist channel \texorpdfstring{$(m = N_f)$}{(m = N_f)}}

The Nyquist edge channel is handled analogously to the DC channel but
centered at the Nyquist frequency $f_{\rm max}$. Its structure
mirrors the DC case with appropriate frequency shifts. In practice,
the Nyquist channel is computed using the same windowed-FFT machinery
as the interior channels, with the frequency shift set to $mN_t/2 = N/2$.
%
\begin{align}\label{eq:wnNf}
  w_{nN_f} &= \sqrt{2}\,\Re\!\left[\sum_{l=N/2-N_t/2}^{N/2-1}
    \tilde{x}[l]\,e^{4\pi i l n/N_t}\,\tilde{\varphi}[l-N/2]\right]
    \nonumber\\
  &\quad + \frac{1}{\sqrt{2}}\,\tilde{x}[N/2]\,\tilde{\varphi}[0]\,.\codelink{https://github.com/pywavelet/wdm_transform/blob/v0.5.0/src/wdm_transform/transforms/xp_backend.py\#L104-L114}
\end{align}

Collecting the results of each of the $w_{nm}$ coefficients above, the forward transform reads
\begin{widetext}
\begin{align}
\label{eq:forward_transform_discrete}
w_{nm} = \sqrt{2}(-1)^{mn}\Re\,
\begin{cases}
\,\displaystyle\sum_{l=1}^{N_t/2 - 1} \tilde{x}[l]\, \tilde{\varphi}[l]\, e^{4\pi i n l/N_t}\,
+ \dfrac{1}{2}\, \tilde{x}[0]\, \tilde{\varphi}[0],
& m = 0, \\[12pt]
 C_{nm}^{*} \displaystyle\sum_{l' = -N_t/2}^{N_t/2 - 1} \tilde{x}[l' + mN_t/2]\, \tilde{\varphi}[l']\, e^{2\pi i l' n/N_t}\,,
& 1 \le m \le N_f - 1, \\[12pt]
\displaystyle\sum_{l = N/2 - N_t/2}^{N/2 - 1} \tilde{x}[l]\, \tilde{\varphi}[l - N/2]\, e^{4\pi i n l/N_t}
+ \dfrac{1}{2}\, \tilde{x}[N/2]\, \tilde{\varphi}[0],
& m = N_f,
\end{cases}
\codelink{https://github.com/pywavelet/wdm_transform/blob/v0.5.0/src/wdm_transform/transforms/xp_backend.py\#L65-L116}
\end{align}
\end{widetext}
holding true for $N_f \in 2\mathbb{N}$. 




\subsection{Backward Transform}

From the orthogonality condition of the Wilson basis, the inverse WDM transform is for the Fourier domain $\tilde{x}$ is immediate 

\begin{align}\label{eq:inverse_fourier}
  \tilde{x}_l = \sum_{n=0}^{N_t-1}\sum_{m=0}^{N_f}\tilde{g}_{nm}[l]\,w_{nm}\,,\codelink{https://github.com/pywavelet/wdm_transform/blob/main/src/wdm_transform/transforms/xp_backend.py\#L366}
\end{align}
found via direct computation
%
\begin{align}
\sum_{n=0}^{N_t-1}\sum_{m=0}^{N_f}\tilde{g}_{nm}[l]\,w_{nm} &=\sum_{n=0}^{N_t-1}\sum_{m=0}^{N_f}\tilde{g}_{nm}[l]\,\sum_{l'=0}^{N-1}\tilde x[l']\tilde g^*_{nm}[l']  \\ \nonumber &=\sum_{l'=0}^{N-1}\tilde x[l']\delta_{ll'}=\tilde x[l]\nonumber 
\end{align}
%
where the first equality uses Eq.~\eqref{wnm_def} and the second using Eq.~\eqref{eq:orthonormality}. 

Specifically, the sum over $n$
produces $N_t\delta_{ll'}$ factors via the discrete orthogonality of
the complex exponentials, and the sum over $m$ covers the full
frequency axis because adjacent windows satisfy
$\tilde{\varphi}^2(f) + \tilde{\varphi}^2(f - \Delta F)
= 1/\Delta F$ in their overlap region (proven in
Appendix~\ref{app:normalization_Meyer}). 
Together these two
identities ensure exact round-trip reconstruction: forward followed
by inverse returns the original data to machine precision, as
verified numerically in Section~\ref{sec:code_package}. We derive explicit formulae for the backward transformations given below \ref{app:inverse_transform}

\begin{widetext}
\begin{align}
\label{eq:inverse_transform_discrete}
\tilde{x}_{m}[l] = \dfrac{1}{\sqrt{2}}
\begin{cases}
\tilde{\varphi}[l]\, \displaystyle\sum_{n=0}^{N_t - 1} w_{n0}\, e^{-4\pi i n l / N_t},
& m = 0,\ \ 0 \le l < N_t/2, \\[10pt]
\tilde{\varphi}[l - mN_t/2]\, \displaystyle\sum_{n=0}^{N_t - 1} C_{nm}\, w_{nm}\, e^{-2\pi i n l / N_t},
& 1 \le m \le N_f - 1,\ \ (m{-}1)\tfrac{N_t}{2} \le l < (m{+}1)\tfrac{N_t}{2}, \\[10pt]
\tilde{\varphi}[l - N/2]\, \displaystyle\sum_{n=0}^{N_t - 1} w_{nN_f}\, e^{-4\pi i n l / N_t},
& m = N_f,\ \ N/2 - N_t/2 \le l \le N/2,
\end{cases}
\codelink{https://github.com/pywavelet/wdm_transform/blob/main/src/wdm_transform/transforms/xp_backend.py\#L119-L176}
\end{align}
\end{widetext}
where the original Fourier domain expression $\tilde{x}[l]$ on the output data is given via
\begin{equation}
\tilde{x}[l] = \sum_{m=0}^{N_f} \tilde{x}^{m}[l]
\end{equation}
Equation \eqref{eq:inverse_transform_discrete} is the first time that the inverse WDM transform has been given in the literature to date. 

\section{Statistics within the WDM formalism}\label{sec:stats}

\subsection{The WDM likelihood}

For a Gaussian frequency-domain noise process characterized by a covariance matrix $\boldsymbol{\tilde \Sigma} = \mathbb{E}[\boldsymbol{\tilde x}\boldsymbol{\tilde x}^T]$,
the log-likelihood of the data $\mathbf{\tilde d}$ given a signal model $\mathbf{\tilde h}$, characterised via residuals $\boldsymbol{\tilde x} = \boldsymbol{\tilde d} - \boldsymbol{\tilde h}$ is \cite{Burke:2025bun}
\begin{equation}\label{eq:frequency_domain_like}
    \log \mathcal{L}(d|\boldsymbol{\theta}) = -\frac{1}{2}\tilde{\boldsymbol{x}}^{\dagger}\tilde{\boldsymbol{\Sigma}}^{-1}\tilde{\boldsymbol{x}} - \ln\det \left(\frac{1}{N}\tilde{\boldsymbol{\Sigma}}\right) - \frac{N}{2}\ln(2\pi)\,.
\end{equation}
If the noise is assumed to be stationary, the frequency domain noise covariance matrix returns a purely diagonal matrix. In this section, we will perform a similar calculation to ~\cite{Burke:2025bun} and transform the frequency domain likelihood \eqref{eq:frequency_domain_like} into the WDM likelihood used for parameter inference.  

We first define the complex rectangular matrices $\boldsymbol{G}\in \mathbb{C}^{N_t(N_f +1) \times N}$ with adjoint $\boldsymbol{G}^{\dagger} \in \mathbb{C}^{N \times N_t(N_f + 1)}$
via the linear maps with action on the frequency domain vector $\tilde{\boldsymbol{x}}\in\mathbb{C}^{N}$ gives $\boldsymbol{w} = \boldsymbol{G}\tilde{\boldsymbol{x}}$ with inverse transform $\tilde{\boldsymbol{x}} = \boldsymbol{G}^{\dagger}\boldsymbol{w}$. In component form, one should interpret $(\boldsymbol{G})_{\alpha,l} = \tilde{g}^{\star}_{nm}[l]$, with components of $(\boldsymbol{G})_{\alpha,l}$ here read as row $\alpha \mapsto (n,m)$ evaluated at frequency index $l$. In component form, we can recall the results of \eqref{eq:orthonormality}:
\begin{align}
    \left(\mathbf{G}^\dagger \mathbf{G}\right)_{ll'} 
    &= \delta_{ll'}, \label{eq:orthogonality_likelihood_proof}\\[6pt]
    \left(\mathbf{G}\mathbf{G}^\dagger\right)_{(n,m),(p,q)}
    &= \delta_{m,q}
    \begin{cases}
        \dfrac{1}{2}\left(\delta_{n,p} + \delta_{n,\,p\pm N_t/2}\right), \\[6pt]
        \delta_{n,p}, 
    \end{cases} \label{eq:orthonormality_like_cases}
\end{align}
%
for the first case in Eq.~\eqref{eq:orthonormality_like_cases} the zeroth $m = 0$ and $m = N_f$ Nyquist bins and second case the interior bins $m \in \{1,\ldots, N_f - 1\}$.
From the orthogonal conditions Eqs. (\ref{eq:orthogonality_likelihood_proof}, \ref{eq:orthonormality_like_cases}), one can see that $\boldsymbol{G}^{\dagger}\boldsymbol{G} = \mathbb{I}_{N} \neq \boldsymbol{G}\boldsymbol{G}^{\dagger}$.

As a consequence the forward map $\boldsymbol{w} = \boldsymbol{G}\tilde{\boldsymbol{x}}$ embeds the signal space $\mathbb{C}^N$ into the larger packed space $\mathbb{C}^{N_t(N_f+1)}$, and the left-inverse $\boldsymbol{G}^\dagger$ recovers the signal,
\begin{equation}
    \tilde{\boldsymbol{x}} \xmapsto{\;\boldsymbol{G}\;} \boldsymbol{w} = \boldsymbol{G}\tilde{\boldsymbol{x}} \xmapsto{\;\boldsymbol{G}^\dagger\;} \boldsymbol{G}^\dagger\boldsymbol{G}\tilde{\boldsymbol{x}} = \tilde{\boldsymbol{x}} \,,
\end{equation}
so the forward-then-inverse round-trip is exact for every signal. The reverse round-trip, by contrast, fails for an arbitrary packed array $\boldsymbol{w}$ that is not itself the transform of a signal:
\begin{equation}
    \boldsymbol{w} \xmapsto{\;\boldsymbol{G}^\dagger\;} \boldsymbol{G}^\dagger\boldsymbol{w} \xmapsto{\;\boldsymbol{G}\;} \boldsymbol{G}\boldsymbol{G}^\dagger\boldsymbol{w} \neq \boldsymbol{w} \,,
\end{equation}
%
The WDM covariance matrix $\boldsymbol{\Sigma}_{w}=\mathbb{E}[\boldsymbol{w}\boldsymbol{w}^{\dagger}]\in\mathbb{R}^{N_t(N_f+1)\times N_t(N_f +1)}$ therefore relates to the one in frequency domain via
%
\begin{align}
    \tilde{\boldsymbol{\Sigma}} &= \mathbb{E}[\tilde{\boldsymbol{x}}\tilde{\boldsymbol{x}}^{\dagger}] = \boldsymbol{G}^{\dagger}\mathbb{E}[\boldsymbol{w}\boldsymbol{w}^{\dagger}]\boldsymbol{G} = \boldsymbol{G}^{\dagger}\boldsymbol{\Sigma}_{w}\boldsymbol{G} \label{def:cov_freq_matrix}\\
    \boldsymbol{\Sigma}_w &= \mathbb{E}[\boldsymbol{w}\boldsymbol{w}^{\dagger}] = \boldsymbol{G}\mathbb{E}[\tilde{\boldsymbol{x}}\tilde{\boldsymbol{x}}^{\dagger}]\boldsymbol{G}^{\dagger} = \boldsymbol{G}\tilde{\boldsymbol{\Sigma}}\boldsymbol{G}^\dagger \label{def:cov_wdm_matrix}
\end{align}
%
for $\boldsymbol{\Sigma}_{w}\in\mathbb{R}^{N_t(N_f+1)\times N_t(N_f +1)}$ the \emph{real valued} WDM noise covariance matrix. 

Now, for stationary and circulant noise the frequency domain covariance matrix $\tilde{\Sigma}$ is trivially invertible. However, we have that $N = \text{Rank}(\boldsymbol{G}^{\dagger}\boldsymbol{G}) = \text{Rank}(\boldsymbol{G}) = \text{Rank}(\boldsymbol{G}^{\dagger})$, and that $\text{Rank}(\boldsymbol{\Sigma}_w)  \leq N$. Since the matrix $\boldsymbol{\Sigma}_{w}$ does not have full rank, i.e., $\text{Rank}(\Sigma_{w}) \neq (N + N_t)$, we deduce that $\Sigma_{w}$ is not invertible over the full span of pixels $n \in \{0,\ldots, N_t -1\}$ and $m\in\{0, \ldots, N_f\}$. 

To make progress, we define the Moore-Penrose pseudo inverse $\boldsymbol{\Sigma}^+_{w}$ of the matrix $\boldsymbol{\Sigma}_w$
%
\begin{equation}\label{eq:pseudo_inv_cov}
\boldsymbol{\Sigma}_w^{+} = \boldsymbol{G}\tilde{\boldsymbol{\Sigma}}^{-1}\boldsymbol{G}^{\dagger}\,,
\end{equation}
%
and we can build the WDM likelihood from \eqref{eq:frequency_domain_like}
%
\begin{align}\label{eq:likelihood_G_Sigma_w}
    \ln\mathcal{L}(w|\boldsymbol{\theta}) &= -\frac{1}{2}\boldsymbol{w}^{T}\boldsymbol{\Sigma}_{w}^{+}\boldsymbol{w}  \\
    & \quad - \ln\det \left(\frac{1}{N_fN_t}\boldsymbol{G}^{\dagger}\boldsymbol{\Sigma}_{w}\boldsymbol{G}\right) - \frac{N_fN_t}{2}\ln(2\pi)\,.
\end{align}


\subsection{The WDM Inner Product}
In a similar way to a frequency domain analysis, under the assumptions in Appendix \ref{app:stats} we obtain the approximated expression for the noise-weighted inner product in the WDM domain 
\begin{align}
    (a|b)_{\mathrm{WDM}} &= \boldsymbol{w_a}^T(\boldsymbol{\Sigma}^{\rm int}_w)^{-1}\boldsymbol{w_b} \\
    & = \sum_{\substack{n \\ m\in\mathrm{int}}} \frac{w_{a,nm}\,w_{b,nm}} {S_{nm}}\,,\label{eq:wdm_inner}\codelink{https://github.com/pywavelet/wdm_transform/blob/main/src/wdm_transform/signal_processing.py\#L52-L60}
\end{align}
%
\gm{Add $S_{nm}$ formula and link it to the appendix}
which reduces to the frequency domain inner product via matrix algebra
%
\begin{align}
    (a|b)_{\mathrm{WDM}} &= \boldsymbol{w_a}^T(\boldsymbol{\Sigma}^{\rm int}_w)^{-1}\boldsymbol{w_b}  \\ 
    &= \tilde{x_a}^{\dagger}\boldsymbol{G}^{\dagger}\boldsymbol{G}(\tilde{\Sigma})^{-1}\boldsymbol{G}\boldsymbol{G}^{\dagger}\boldsymbol{\tilde x_b} \\
    & = \tilde{\boldsymbol{x_a}}^{\dagger}(\tilde{\boldsymbol{\Sigma}})^{-1}\tilde{\boldsymbol{x_b}} \\
    & = \frac{4N}{\Delta t} \text{Re}\left(\sum_{l=-N/2}^{N/2 - 1}\frac{\tilde{x_a}^{\star}[l]\tilde{x_b}[l]}{S_{n}[l]}\right) \\
    & =(a|b)_{\text{freq}}\,.
\end{align}
Where each vector and Wilson basis matrix $\boldsymbol{G}$ only contains full time information $n\in \{0,\ldots,N_t-1\}$ but, as a first working approximation, we exclude the zeroth/Nyquist bin $m\in\{1,\ldots,N_f -1\}$\,. 
Similarly, the SNR, Fisher Matrix and bias \cite{Cutler:2007mi} estimates read
\begin{align}
    \rho^{\text{opt}} &= (h|h)_{\rm WDM}\,, \qquad \text{SNR} \\
    \rho(h) &= \frac{(d \,|\, h)_{\rm WDM}}{\sqrt{(h \,|\, h)_{\rm WDM}}}\,, \quad \text{Effective SNR}\\
    % \codelink{https://github.com/pywavelet/wdm_transform/blob/main/src/wdm_transform/signal_processing.py\#L52-L60}
    \Gamma_{ij} &= (\partial_{i}h|\partial_{j}h)_{\rm WDM}\,, \quad \text{Fisher Matrix} \label{eq:wdm_fisher} \\
    \Delta \boldsymbol{\theta}^{i} &= (\Gamma^{-1})^{ij}(\partial_{i}h|n + \delta h)_{\rm WDM}\,. \quad \text{Cutler Valisneri}
\end{align}
%
These expressions assume that the diagonal
approximation~\eqref{eq:cov_interior_WDM} holds. When it does not ---
for example, when the noise PSD varies rapidly compared to $\Delta
T$, or when the noise is strongly non-Gaussian --- the off-diagonal
covariance terms must be retained and the simple pixel-wise sums
above become inadequate. Figure~\ref{fig:wdm_orthogonality_decorrelation}
provides a practical
diagnostic: if the off-diagonal entries of the empirical covariance
remain small after whitening, the diagonal
likelihood~\eqref{eq:wdm_whittle} is a safe approximation. A natural question comes to light -- when does the diagonal approximation hold? 

The WDM tiling obeys $\Delta T\,\Delta F = 1/2$ (see
Eq.~\eqref{DTDF}), so each atom occupies the minimum allowed
time-frequency area consistent with the uncertainty principle, though
the actual time-bandwidth product of the cosine-tapered Meyer window
exceeds this minimum. The diagonal covariance~\eqref{eq:cov_interior_WDM}
holds whenever the noise PSD varies slowly compared to this joint
support. Concretely, this is the condition of \emph{local
stationarity}: spectral variations on timescales shorter than $\Delta
T$, or spectral features narrower than $\Delta F$, will induce
off-diagonal covariance terms. The WDM tiling therefore poses an
inherent trade-off: choosing finer frequency resolution (larger
$\Delta T$) suppresses spectral leakage within each atom but degrades
temporal tracking of non-stationarity. In practice a single choice of
$N_f$ often suffices for mildly non-stationary noise, but data with
sharp spectral features or rapid temporal variation may require a
multi-resolution strategy, as discussed by
Cornish~\cite{Cornish:2020odn}.



\begin{figure*}[h]
  \centering
  \includegraphics[width=1\linewidth]{figures/wdm_orthogonality_decorrelation.pdf}
  \caption{
  \textbf{Orthogonality structure and empirical decorrelation of the
sampled WDM transform.} 
All panels use the small tiling $N_t = N_f = 8$,
$N = N_t N_f = 64$, with packed atoms flattened by $\alpha = n(N_f + 1) + m$.
Panels (a), (c), and (d) live in the $N_t(N_f+1) = 72$-dimensional packed
coefficient space, while panel (b) lives in the $N$-dimensional signal
space. The asymmetry $GG^\dagger \in \mathbb{C}^{72\times 72}$ versus
$G^\dagger G \in \mathbb{C}^{64\times 64}$ is intrinsic to the rectangular
embedding $G : \mathbb{C}^N \to \mathbb{C}^{N_t(N_f+1)}$ and is the
matrix-level expression of the fact that $G^\dagger G = \mathbb{I}_N \neq
GG^\dagger$.
Faint grid lines in the panels mark the block boundaries $\alpha = n(N_f+1)$
for $n = 0, 1, \ldots, N_t$, separating successive time slices and making
the $\pm N_t/2$ offset structure easy to read off by eye.
\textbf{(a)} The packed-space Gram matrix
$(GG^\dagger)_{(n,m),(p,q)} = \sum_\ell g_{nm}[\ell]\,g^*_{pq}[\ell]$. Interior
channels ($m = 1,\ldots,N_f-1$) form an identity block, while the DC and
Nyquist edge channels ($m \in \{0, N_f\}$) carry the off-diagonal entries at
$p = n \pm N_t/2$ predicted by the first case of
Eq.~\eqref{eq:orthonormality_like_cases}. These off-diagonals span the
$N_t$-dimensional redundant subspace by which $GG^\dagger$ fails to be the
identity. \textbf{(b)} The signal-space product
$(G^\dagger G)_{\ell\ell'} = \sum_{n,m} g_{nm}[\ell]\,g^*_{nm}[\ell']$,
indexed by sample indices $(\ell,\ell')$ rather than the packed pair, equals
$\mathbb{I}_N$ to machine precision ($\max|G^\dagger G - \mathbb{I}| \sim
10^{-15}$), demonstrating that the forward-then-inverse round trip is exact.
\textbf{(c, d)} Empirical correlation matrices of the packed WDM coefficients
estimated from $4000$ realisations of white \textbf{(c)} and coloured
\textbf{(d)} stationary Gaussian noise. The coloured PSD is
$S(f) \propto [1 + (f/f_c)^4]^{-1}$ with $f_c = 0.06$ (unit sampling rate).
The interior bulk is nearly decorrelated for white noise and develops
visible band structure for coloured noise. 
The exact $\pm 1$ entries at the
same $\pm N_t/2$ offsets seen in panel (a) are the empirical signature of
the edge-channel redundancy: $w_{n,0}$ and $w_{n + N_t/2,\,0}$ (and likewise
for $m = N_f$) are linearly dependent copies of the same underlying real
degree of freedom, so their Pearson correlation is identically $\pm 1$ for
any input. The reported $\text{mean}\,|\rho_{\rm off}|$ is computed over
interior pairs only, so it summarises genuine decorrelation rather than
storage redundancy.}
  \label{fig:wdm_orthogonality_decorrelation}
  \script{wdm_orthogonality_decorrelation.py}
\end{figure*}


\section{\package\ Package}
\label{sec:code_package}

\ob{I also think a sexy flowchart describing the package would be amazing here. At least with all of the various modules. The main thing we are selling here is (1) an understanding of the mathematical formalism (done) and (2) the code base. A simple flow chart so that any person reading the paper in 15 seconds can stop on would be really useful. At least then they can see: forward/backwards transforms, docs, unittests and the important modules.}
\avi{sure can do..}

The \package\ package provides NumPy and JAX backends for the
Wilson--Daubechies--Meyer (WDM) transform, both implementing the
packed $(N_t, N_f+1)$ coefficient layout and forward/inverse
conventions described above.  The primary interface exposes the full
forward and inverse transforms; in addition, band-limited routines
allow the transform to be evaluated over a specified subset of
frequency channels, reducing compute cost when only a portion of the
time-frequency plane is needed.

From the implementation point of view, the dominant operations in the
forward transform are one length-$N$ FFT of the input data together
with one length-$N_t$ sub-band transform for each stored channel. This
gives a cost of
\begin{align}
  T_{\rm forward} = \mathcal{O}(N \log N) + \mathcal{O}(N \log N_t),
\end{align}
which is asymptotically $\mathcal{O}(N \log N)$ because $N=N_tN_f$
and $N_t \leq N$ 
%\ob{Actually, hang on a minute. The FT only scales as $\mathcal{O}(N\log N)$ if and only if The FFT computations are acting on powers of two or $N$ interestingly, has a simple prime factorisation primes $N = 2^{n_1}\cdot 3^{n_2} \cdots$. I.e., if you set $N = 2 \times \text{large prime}$, then you'll still get $\mathcal{O}(N \log N)$ but a noticeable slow down due to the massive constant at the front. In reality, this probably doesn't matter. It would be hilarious of $N = 2 \times \text{huge prime}$, but we should argue ``asymptotically" in a tighter sense. I can do this if you wish?} \gm{I take care of it, after Avi adds to figure 3}. 
The inverse transform has the same scaling up to
constant factors: it reconstructs the full spectrum from the packed
channels and then applies one inverse FFT back to the time domain. In
both directions, the stored coefficient grid contains $\mathcal{O}(N)$
real numbers, and the temporary working arrays used by the
implementation are also linear in $N$, so the overall memory footprint
is $\mathcal{O}(N)$.

We validate the implementation using round-trip reconstruction tests.
For an input time series $x$, we compute its packed WDM coefficients,
apply the inverse transform, and compare the reconstructed series
$\hat{x}$ with the original. Across the benchmarked problem sizes and
backends, the relative reconstruction error
\begin{align}
  \epsilon_{\rm rt} = \frac{\|\hat{x}-x\|_2}{\|x\|_2}
\end{align}
remains at the level of double-precision floating-point roundoff,
$\epsilon_{\rm rt} \lesssim 10^{-16}$, providing a direct numerical
check that the packed layout, including the DC and Nyquist edge
channels, is implemented consistently across both backends.

All timings are drawn from the benchmark suite
included in the package documentation. Each data point reports the
median wall-clock time over 7 repeated runs in 64-bit floating-point
arithmetic; JAX JIT compilation is excluded via a single warmup
evaluation, and GPU timings assume arrays already resident on the
device (host--device transfer is not included). Benchmarks were
performed on Google Colab using a CPU node (2$\times$ Intel Xeon
vCPU, 2.0--2.2\,GHz, 12.7\,GB RAM) and a GPU node (NVIDIA T4,
16\,GB VRAM).

Figure~\ref{fig:runtime_forward} shows the forward-kernel runtime as a
function of input size $N$ for the NumPy (CPU), JAX (CPU), and JAX
(GPU) backends, together with the batched-execution speedup
$t_\mathrm{serial}/t_\mathrm{batch}$. Here, $t_\mathrm{serial}$ is
the wall-clock time for three independent single-channel calls and
$t_\mathrm{batch}$ is the time for a single call that processes all
three channels simultaneously.

\begin{figure}[tbp]
  \centering
  \includegraphics[width=1\linewidth]{figures/runtimes.pdf}
  \caption{\textbf{Forward WDM runtime versus input length.}
    \emph{Top}: median wall-clock runtime over 7 runs after one warmup
    call; shaded bands span the 16th--84th percentile interval.  The
    black dashed line shows the $N\log_2 N$ reference scaling.
    \emph{Bottom}: batched-execution speedup
    $t_\mathrm{serial}/t_\mathrm{batch}$, where the serial baseline
    applies the kernel independently to each of three channels and the
    batched call processes all three simultaneously.  Reported timings
    exclude JAX JIT compilation and host--device data transfer.
    \ls{I think that calling the y-axis of the bottom plot just "speed-up" would make it easier for the reader}
    \avi{sure, can do. Also i will plot nlog(nt) and FFT timing, up to $N=10^7$.}
    }
  \label{fig:runtime_forward}
  \script{benchmark_runtime_snapshot.py}
\end{figure}

At large $N$ all three backends approach $\mathcal{O}(N\log N)$
scaling, consistent with the dominant length-$N$ FFT in the forward
transform.  Deviations at small $N$ reflect different overhead
regimes: per-call Python and memory-allocation costs for NumPy, and
kernel-dispatch latency for JAX.  The JAX GPU runtime is nearly flat
below $N \approx 10^5$, indicating that transforms at those sizes are
dominated by kernel-launch overhead rather than arithmetic throughput;
beyond this threshold GPU efficiency improves and, at $N = 10^6$, the
GPU scalar runtime (${\approx}\,13\,\mathrm{ms}$) is roughly
$35\times$ lower than NumPy (${\approx}\,460\,\mathrm{ms}$).

The speedup panel shows that batched execution yields no measurable
benefit for NumPy ($t_\mathrm{serial}/t_\mathrm{batch} \approx 1$
throughout), as NumPy does not exploit the extra batch dimension for
additional parallelism.  JAX (CPU) recovers $1.5$--$3\times$ via
JIT-compiled vectorisation over the batch axis.  JAX (GPU) achieves
up to ${\sim}\,4\times$ at small $N$, where all three channels fit
simultaneously into the device pipeline; the speedup decreases to
${\approx}\,1.5\times$ at $N = 10^6$ as individual transforms
approach device saturation.  The $\epsilon_{\rm rt} \lesssim 10^{-16}$
round-trip errors established above hold for all backends and sizes;
the runtime differences therefore reflect implementation efficiency
rather than any loss of numerical fidelity.

These benchmarks characterise the present NumPy and JAX implementations
under the conventions used in this paper, rather than providing an
exhaustive comparison with existing pipeline-specific WDM codes.  The
main conclusion is that the JAX backend can substantially reduce the
cost of large WDM transforms while preserving the same floating-point
reconstruction accuracy as the NumPy implementation.

\subsection{Resolved galactic-binary inference in a galactic background}
\label{sec:lisa_gb_example}

\gm{After discussing with Ollie, we think it is better to do the inference with just instrumental noise, which is stationary. Avi, what do you think about it? If you agree, could you please re-run the analysis? We can take care of the rewriting of the section (cutting off the galactic background formulas)}
\ls{consider looking into \cite{Du:2025fes} for GBs in non stationary noise ;)}
\gm{Our concern is just for the sake of clarity. These days Ollie and I are working on nonstationary PSD, but even if we have solid results I'd rather put them in a follow-up. This work should be as clean and simple as possible, what do you all think?}
As an end-to-end inference example, we generated a one-year LISA-like
dataset containing stationary instrumental noise, a stochastic galactic
foreground, and one resolved galactic binary (GB). Both the
frequency-domain and WDM-domain likelihoods use the same stationary
noise model throughout; this is an intentional controlled comparison
designed to isolate the effect of the WDM basis change from any
modelling differences. The effective PSD $S^{\mathcal{I}}_n(f)$ in each channel $\mathcal{I}$
includes both the instrumental noise contribution and the
time-averaged galactic foreground. Figure~\ref{fig:lisa_gb_data} shows
the characteristic-strain comparison between the injected source and
the channel noise model, and
Figure~\ref{fig:lisa_gb_corner} shows the overlaid posterior
constraints from the two analyses.

The data set is constructed on a frequency grid $\in 
(10^{-4}$, $3\times 10^{-3}\,\mathrm{Hz})$. The unresolved galactic
foreground is described by the phenomenological PSD \cite{Karnesis:2021tsh}
%
\begin{align}
  S_{\rm gal}(f) &=
  A_{\rm gal} f^{-7/3}
  \exp\!\left[-\left(\frac{f}{f_{1,{\rm gal}}}\right)^{\alpha_{\rm gal}}\right]
  \nonumber\\
  &\quad\times
  \left[1 + \tanh\!\left(\frac{f_{{\rm knee},{\rm gal}} - f}{f_{2,{\rm gal}}}\right)\right],
\end{align}
%
using a standard smooth fit for the LISA galactic foreground. To
project this foreground into the detector channels we combine that
frequency model with an anisotropic Milky Way sky map obtained by
line-of-sight integration through the bulge-plus-disk density model \cite{Mentasti:2023uyi}
%
\begin{align}
  \rho_{\rm gal}(x,y,z) =
  \rho_0\left[
    A e^{-r^2/R_b^2}
    + (1-A)e^{-u/R_d}\,{\rm sech}^2(z/Z_d)
  \right],
\end{align}
%
where $r^2=x^2+y^2+z^2$ and $u^2=x^2+y^2$. We evaluate the
time-dependent detector response using equal-arm LISA orbits~\cite{lisaorbits}
sampled at $90$ orbital times, and use that response to map the
anisotropic foreground into the A,E,T TDI channels, assuming that the
spectrum is stationary within each of these $90$ orbital segments. We
then draw complex Gaussian realisations of both the instrumental noise
and the chunked galactic foreground and stitch them into a full-year
spectrum. The resolved source is injected afterward as a TDI~1.5 A,E,T
waveform computed with \texttt{jaxgb}~\cite{jaxgb}, using equal-arm
LISA orbits and local Fourier support $n=256$. In the realization shown here the cadence is $\Delta t = 166.7\,\mathrm{s}$, $T_{\rm obs}=1\,\mathrm{yr}$, and the injected binary has $f_0 = 1.3862897\,\mathrm{mHz}$, $\dot f = 1.32\times 10^{-18}\,\mathrm{Hz\,s^{-1}}$, $A = 1.19\times 10^{-23}$, and $\phi_0 = 1.442\,\mathrm{rad}$.

\gm{I will define sigma gal, and check everything here, after Ollie is ready with the math of the statistics}\ob{I think we are ready for you to proceed here. }

The resolved GB model uses the standard eight source parameters
$(f_0,\dot f,A,\alpha,\delta,\psi,\iota,\phi_0)$, but for this
controlled comparison we sample only
$(f_0,\dot f,A,\phi_0)$ and hold sky position and orientation fixed to
their injected values. The carrier frequency is drawn near the fixed
reference value $f_{0,\rm ref}=1.38591\,\mathrm{mHz}$ through
$\delta\!\ln f_0 \sim \mathcal{U}[-10^{-3},10^{-3}]$, while
$\dot f$ and $A$ are drawn from truncated Gaussians in log space over
$[5\times 10^{-19},\,4\times 10^{-18}]\,\mathrm{Hz\,s^{-1}}$ and
$[6\times 10^{-24},\,1.7\times 10^{-23}]$, respectively.

The frequency-domain analysis uses a local A/E/T band centered on the
prior support of the source and evaluates the Whittle likelihood
% \begin{align}
%   \ln \mathcal{L}_{\rm freq}(\theta)
%   &\propto
%   -\sum_{c\in\{A,E,T\}} \sum_k
%   \left[
%     \ln S_c(f_k)
%     + \frac{2\Delta f\,\bigl|\Delta t\,[d_{c,k}-h_{c,k}(\theta)]\bigr|^2}{S_c(f_k)}
%   \right],
% \end{align}
\begin{equation}
\begin{split}
  \ln \mathcal{L}_{\rm freq}(\boldsymbol{\theta}) \propto {} & -\sum_{\mathcal{I}\in\{A,E,T\}} \sum_k \bigg[ \ln S^{(\mathcal{I})}_n[k] \\
  & + \frac{2N\,\bigl|\tilde{d}^{(\mathcal{I})}[k] - \tilde{h}^{(\mathcal{I})}[k;\boldsymbol{\theta}]\bigr|^2}{\Delta tS^{(\mathcal{I})}_n[k]} \bigg],
\end{split}
\end{equation}
where for each channel $\mathcal{I}$ and frequency bin $k$, $\tilde{d}^{(\mathcal{I})}[k]$ and $\tilde{h}^{(\mathcal{I})}[k;\boldsymbol{\theta}]$ are the complex local Fourier modes and
$S^{(\mathcal{I})}_n[k]$ is the stationary effective PSD, namely instrumental noise\footnote{We adopt the Gaussian two-parameter model for the instrumental noise PSD~\cite{LISACosmologyWorkingGroup:2022kbp}, using Eqs.~(B14) and~(B15) therein.}
plus the time-averaged galactic
foreground contribution in channel $c$. The WDM analysis starts from
the same local Fourier patch, embeds the template into that
patch, and applies the WDM transform with
$N_t=32$, $a=1/3$, and $d=1$. The resulting WDM likelihood is the
diagonal Gaussian specialization of Eq.~\eqref{eq:wdm_whittle},
\begin{align}
  \ln \mathcal{L}_{\rm WDM}(\theta)
  &= -\frac{1}{2}\sum_{c\in\{A,E,T\}} \sum_{n,m}
  \Bigg[
    \ln(2\pi \sigma^2_{c,m})
    \nonumber\\
  &\qquad\qquad
    + \frac{\left(w^{(d)}_{c,nm}-w^{(h)}_{c,nm}(\theta)\right)^2}{\sigma^2_{c,m}}
  \Bigg],
\end{align}
with a stationary per-pixel variance following from
Eq.~\eqref{eq:snm_stationary} with $\Delta F = f_{\rm Nyq}/N_f$ over the interior $m \in \{1,\ldots,N_f\}$
\begin{align}
  \sigma^2_{c,m} = S^{\text{WDM}}_{nm} = \delta_{np}\delta_{pq}\frac{ N S_c(f_m)}{2\Delta t},
\end{align}
broadcast across the WDM time bins. In other words, this example uses
the same stationary PSD in both domains and asks whether the WDM
transform reproduces the frequency-domain posterior under that shared
noise model.
\avi{Leave a note here on }

Both analyses use the NumPyro~\cite{phan2019composable} implementation
of the No-U-Turn Sampler with a dense mass matrix, target acceptance
probability $0.85$, and maximum tree depth $10$.
Each run uses $800$ warmup steps followed by
$1000$ retained samples from a single chain, initialized near the
injected parameters with small jitters in $\log \dot f$ and $\log A$.
In the realization shown here, both samplers produced zero divergences. The
latent effective sample sizes are $613$--$637$ for the
frequency-domain chain and $590$--$714$ for the WDM chain, while the
corresponding $\hat R$ values lie in the ranges
$0.999$--$1.007$ and $0.999$--$1.016$. The truth-band optimal SNRs are
$22.5469$ in the frequency-domain calculation and $22.5462$ in the
WDM calculation. All reported $90\%$ credible intervals contain the
injected truth in both runs, and the posterior medians differ only by
$9.0\times 10^{-11}\,\mathrm{Hz}$ in $f_0$,
$2.2\times 10^{-20}\,\mathrm{Hz\,s^{-1}}$ in $\dot f$,
$5.2\times 10^{-26}$ in $A$,
$6.1\times 10^{-3}\,\mathrm{rad}$ in $\phi_0$, and
$0.10$ in SNR. Within the numerical precision of this study, the WDM
and frequency-domain analyses are therefore statistically
indistinguishable for this seeded resolved-GB example.

\begin{figure}[t]
  \centering
  \includegraphics[width=\linewidth]{figures/lisa_gb_demo_data.pdf}
  \caption{\textbf{LISA galactic background: Strain comparison.} Characteristic-strain comparison between the injected resolved galactic binary and the channel noise model including the galactic foreground.\gm{Quentin says it's a bit confusing to have A data and PSD in the same plot (A data not following the total PSD curve because it is not a PSD, but $|h|^2$). Shall we renormalize the data to match the PSD?}}
  \label{fig:lisa_gb_data}
  \script{lisa_gb_data.py}
\end{figure}

\begin{figure}[t]
  \centering
  \includegraphics[width=\linewidth]{figures/lisa_gb_demo_corner.pdf}
  \caption{\textbf{LISA galactic background: Posterior analysis.} Overlaid posterior corner plot for the WDM-domain and frequency-domain analyses of the same realization. The two posterior families overlap closely in $(f_0, \dot{f}, A, \phi_0)$ and in the derived optimal SNR.}
  \label{fig:lisa_gb_corner}
  \script{lisa_gb_corner.py}
\end{figure}


To test whether the agreement seen in Figures~\ref{fig:lisa_gb_data}
and~\ref{fig:lisa_gb_corner} is robust across a range of signals and
noise draws, we repeated the inference for 100 independent seeds. For
each seed the source parameters $(f_0, \dot{f}, A, \phi_0)$ and the
sky position and orientation were drawn fresh from the prior, and an
independent realisation of the instrumental noise and galactic
foreground was generated; sky position and orientation were then held
fixed at their injected values during sampling. Both likelihoods use
the same time-averaged effective PSD throughout, so this study
isolates the effect of the WDM basis change from any difference in
noise modelling: the non-stationarity of the galactic foreground
enters both analyses as the same shared systematic and does not
contribute to the JSD. For each seed we computed the Jensen--Shannon
divergence (JSD) between the WDM-domain and frequency-domain 1D
marginal posteriors. The JSD is a symmetric, bounded measure of
distributional similarity, with $\mathrm{JSD} = 0$ for identical
distributions. Table~\ref{tab:jsd} summarises the resulting
distributions across the 100 seeds: median values lie at or below
$0.016\,\mathrm{nat}$ for all four parameters, and the
90th-percentile upper bound reaches at most $0.07\,\mathrm{nat}$.
These results confirm that the WDM-domain likelihood faithfully
reproduces the frequency-domain posterior under a shared stationary
noise model, across a representative population of signal parameters
and noise realisations.

% Design notes:
%   - JSD between WDM and frequency-domain 1D marginal posteriors
%   - n = 100 seeds, same signal injection, independent noise draws
%   - Median + [q05, q95] inter-seed interval
%   - SNR omitted (derived quantity; tracks A and f0)
%
\begin{table}[tbp]
  \centering
  \caption{\textbf{Jensen--Shannon divergence between WDM and
    frequency-domain marginal posteriors across 100 independent noise
    realisations.} The injected signal parameters are held fixed; only
    the noise draw varies. JSD is computed on the 1D marginal for each
    sampled parameter. All values in nats; the $[\mathrm{q}_{05},
    \mathrm{q}_{95}]$ interval characterises seed-to-seed variability.}
  \label{tab:jsd}
  \begin{ruledtabular}
    \begin{tabular}{lcc}
      Parameter
        & Median
        & $[\mathrm{q}_{05},\,\mathrm{q}_{95}]$ \\
      \hline
      $f_0$ [Hz]
        & $0.0139$
        & $[0.0086,\;0.0469]$ \\
      $\dot{f}$ [Hz/s]
        & $0.0106$
        & $[0.0069,\;0.0139]$ \\
      $A$
        & $0.0125$
        & $[0.0087,\;0.0227]$ \\
      $\phi_0$ [rad]
        & $0.0153$
        & $[0.0051,\;0.0686]$ \\
    \end{tabular}
  \end{ruledtabular}
\end{table}



The \package\ package enables both a straightforward NumPy-based
development path and high-performance accelerator execution via JAX.
For small data lengths ($N \lesssim 2^{20}$), the NumPy version is
competitive or even faster due to its extremely low overhead.
However, once $N$ exceeds $\sim 2^{20}$, the GPU-accelerated versions
achieve markedly lower runtimes while both backends maintain
round-trip reconstruction errors at floating-point roundoff.
This combination of flexibility, accuracy, and performance makes
\package\ a useful tool for WDM time-frequency analysis of
long-duration gravitational-wave signals.

\section{Conclusions}\label{sec:conclusions}

We have presented a self-contained, pedagogical account of the
Wilson--Daubechies--Meyer (WDM) wavelet-packet transform, building on
the theoretical foundations established in prior
work~\cite{Necula_2012,Klimenko:2005xv,Klimenko:2015ypf,Cornish:2020odn},
and its implementation in the open-source \package\ package.  The
primary goal of this work is to make the discrete, sampled WDM
implementation explicit and reproducible for practitioners new to the
transform, by walking through the full three-case basis definition
(DC, interior, and Nyquist channels), the packed $(N_t, N_f{+}1)$
coefficient layout, the forward and inverse transforms with detailed
derivations, and the conditions under which WDM coefficients may be
treated as approximately independent.

On the implementation side, the NumPy and JAX backends both achieve
exact round-trip reconstruction to floating-point precision. The JAX
backend provides substantially lower runtimes for large transforms,
making it practical to incorporate WDM-domain operations into
inference pipelines that require millions of likelihood evaluations.
\ob{One point to make here -- is the WDM transformation as cheap as a simple FFT? If so, then this makes incorporating the WDM methodology into any FD based analysis will not improve computational costs, implying that there is no costs to using WDM over current FD analyses? Does this make sense?}
\avi{not quite as cheap as a single FFT i think... the forward WDM is an FFT plus $N_f$ windowed inverse FFTs of length $N_t$, so the scaling is $O(N\log N_t)$ rather than $O(N\log N)$. in practice the constant is small enough that the cost is comparable to a few FFTs... maybe we can include FFT timings in the runtime plot to compare.}

Several limitations of the current work should be noted. First, the
package implements only the $d=1$ Meyer window; higher-order windows
($d > 1$) offer better spectral leakage suppression and may be
preferable for data with sharp spectral features.
Second, the paper
does not address the generation of gravitational-wave signals directly
in the WDM domain — where much of the computational advantage of the
WDM representation lies~\cite{Cornish:2020odn} — and we defer fast
WDM-domain waveform generation to future work. Third, the
treatment of data gaps and edge effects, while conceptually
straightforward, requires careful handling of the boundary basis elements
and is left to downstream analyses built on top of the package.%

Fourth, the LISA galactic-binary
inference example of Section~\ref{sec:lisa_gb_example} uses a
time-averaged effective PSD in both the frequency-domain and
WDM-domain likelihoods; this controlled comparison isolates the WDM
basis change but does not exercise the principal advantage of the WDM
representation, namely the ability to assign a time-varying $S_{nm}$
to each pixel and thereby accommodate genuinely non-stationary noise
without off-diagonal covariance terms. Demonstrating this advantage on
realistic LISA non-stationary backgrounds is an important direction for
future work.

Looking ahead, the most immediate extensions are the incorporation of
fast WDM-domain waveform generation for compact binary signals,
the implementation of evolutionary power spectral density estimation
directly on the WDM grid, and integration of the package into
existing LISA data analysis frameworks.  The combination of a
pedagogical reference, a tested open-source implementation, and
GPU-accelerated performance should lower the barrier to adopting the
WDM transform within modern, autodiff-compatible Python inference
pipelines for the long-duration, non-stationary signals expected from
future detectors such as LISA. Because the JAX backend is fully
differentiable, gradient-based samplers (HMC, NUTS) and variational
inference are available out of the box on top of the WDM likelihood,
without any additional differentiation work by the user. Future work should focus on determining the optimal tiling of the WDM space for each source class. A systematic comparison with the short-time Fourier transform, assessing respective advantages and limitations, would help identify the most appropriate time–frequency representation and tiling for different types of sources. 

\section*{Data and Software Availability}
The software developed for this research is available on GitHub at
\url{https://github.com/pywavelet/wdm_transform}. Documentation and
worked examples are available in the repository documentation at
\url{https://github.com/pywavelet/wdm_transform/tree/main/docs}.
The online documentation includes minimal examples demonstrating the
construction of a WDM tiling, application of the forward transform, and
round-trip reconstruction with the inverse transform.

\begin{acknowledgments}
GM acknowledges support from the ANR grant AAPG2024 PRC - GalaxyFIT. O.~Burke acknowledges financial support from the Grant UKRI972 awarded via the UK Space Agency. O.~Burke appreciates insightful discussions with Christopher Moore and Sylvain Marsat during the preparation of this work.  
L. Speri acknowledges support through the European Space Agency (ESA) Research Fellowship in Space Science.
\end{acknowledgments}

\appendix
\onecolumngrid

\section{Normality condition of the $g_{nm}$ basis}\label{app:normality_condition}
%
To remind the reader, We introduce the basis $\tilde g_{nm}(f)$, given by eq.\eqref{eq:gnm_full_discrete}, 
%
\begin{align}
\label{eq:gnm_discr}
\tilde{g}_{nm}[l] =
\dfrac{1}{\sqrt{2}}\begin{cases}
e^{-4\pi i n l/N_t}\, \tilde{\varphi}[l], & m=0, \\[6pt]
e^{-2\pi i n l/N_t}\Big(C_{nm}\,\tilde{\varphi}[l - mN_t/2] + C_{nm}^*\,\tilde{\varphi}[l + mN_t/2]\Big), & 0 < m < N_f, \\[6pt]
e^{-4\pi i n l/N_t}\Big(\tilde{\varphi}[l - N/2] + \tilde{\varphi}[l + N/2]\Big), & m = N_f,
\end{cases}
\end{align}
%
The closed-form expressions below are written
for $d=1$ to keep the algebra explicit. We do remark, however, that the orthonormality and
normalisation results extend to arbitrary $d$ using the incomplete-beta
identity $\nu_d(1-y)+\nu_d(y)=1$. We restrict to
$d=1$ here purely for pedagogical clarity:
%
\begin{align}\label{eq:phi_discrete}
\tilde\varphi[l]&=\sqrt{\frac{2}{N_t}}\begin{cases}
1 & |l_r|<A \\
\cos \left[\frac{\pi}{2}\,\left(
\frac{|l_r|-A}{B}\right)\right] & A \leq|l_r| < A+B \\
0 & |l_r|\geq A+B
\end{cases}\nonumber\\
l_r&=l\frac{2}{N_t}
\end{align}
%
where $l\in\{-\frac{N_t}{2},-\frac{N_t}{2}+1,...,\frac{N_t}{2}-1\}$. Now we verify that
%
\begin{align}
\sum_{nm}&\tilde{g}_{nm}[l]\tilde{g}_{nm}^*[l']=\frac{1}{2}\sum_n e^{-4\pi i n (l-l')/N_t}\, \tilde{\varphi}[l]\,\tilde{\varphi}[l']\nonumber\\
&+\frac{1}{2}\sum_{0<m<N_f}\sum_{n}e^{-2\pi i n (l-l')/N_t}\,\Big(C_{nm}\,\tilde{\varphi}[l - mN_t/2] + C_{nm}^*\,\tilde{\varphi}[l + mN_t/2]\Big)\Big(C_{nm}^*\,\tilde{\varphi}[l' - mN_t/2] + C_{nm}\,\tilde{\varphi}[l' + mN_t/2]\Big)\nonumber\\
&+\frac{1}{2}\sum_n e^{-4\pi i n (l-l')/N_t}\Big(\tilde{\varphi}[l - N/2] + \tilde{\varphi}[l + N/2]\Big)\Big(\tilde{\varphi}[l' - N/2] + \tilde{\varphi}[l' + N/2]\Big)
\end{align}
%
Now we simplify the $0<m<N_f$ term, by noticing that $C_{nm}^2=C_{nm}^{*2}=(-1)^{m}(-1)^{n}$ and $|C_{nm}|^2=1$
%
\begin{align}
\frac{1}{2}\sum_{0<m<N_f}\sum_{n}&e^{-2\pi i n (l-l')/N_t}\,\Big(\tilde{\varphi}[l - mN_t/2]\tilde{\varphi}[l' - mN_t/2] + \tilde{\varphi}[l + mN_t/2]\tilde{\varphi}[l' + mN_t/2]\nonumber\\
&+ (-1)^{m}(-1)^{n}\tilde{\varphi}[l - mN_t/2]\tilde{\varphi}[l' + mN_t/2]+ (-1)^{m}(-1)^{n}\tilde{\varphi}[l + mN_t/2]\tilde{\varphi}[l' - mN_t/2]\Big)
\end{align}
%
one notices that the summation over $n$ can be directly performed, and it vanishes for every value but $l=l'$. In particular, the second line (the mixed terms) always vanishes. A similar result holds for the $m=0,N_f$ cases, therefore we can write
%
\begin{align}
\sum_{nm}&\tilde{g}_{nm}[l]\tilde{g}_{nm}^*[l']=\delta_{ll'}\frac{N_t}{2}\left(\tilde\varphi^2[l]+\sum_{0<m<N_f}\tilde\varphi^2[l-mN_t/2]+\sum_{0<m<N_f}\tilde\varphi^2[l+mN_t/2]+\tilde\varphi^2[l-N/2]+\tilde\varphi^2[l+N/2]\right)
\end{align}
%
Now the thing to notice is that each value of $l$ lays in two and only two of these windows. Let be $\bar m$ and $\bar m+1$ the indices of these two windows, such that $\frac{\bar m N_t}{2}<|l|<\frac{(\bar m+1)N_t}{2}$ ($\bar m$ could be 0, or $\bar m+1$ can be $N_f$). Then the sum becomes
%
\begin{align}
\sum_{nm}&\tilde{g}_{nm}[l]\tilde{g}_{nm}^*[l']=\delta_{ll'}\frac{N_t}{2}\left(\tilde\varphi^2[l-\bar m N_t/2]+\tilde\varphi^2[l-(\bar m+1)N_t/2]\right)
\end{align}
%

Let $l\geq 0$ for sake of simplicity ($l<0$ cases are analogous). Then either $l<\frac{(m+A)N_t}{2}$ or $l\geq\frac{(m+A)N_t}{2}$.
In the first case, we have $\tilde\varphi^2[l-\bar m N_t/2]=\frac{2}{N_t}$ and $\tilde\varphi^2[l-(\bar m+1)N_t/2]=0$, since for $|\frac{2l}{N_t}-(\bar m+1)|>A+B$. In the second case, we have
%
\begin{align}
\sum_{nm}&\tilde{g}_{nm}[l]\tilde{g}_{nm}^*[l']=\delta_{ll'}\left(\cos^2\left[\frac{\pi}{2}\frac{\frac{2l}{N_t}-\bar m -A}{B}\right]+\cos^2\left[\frac{\pi}{2}\frac{\bar m + 1 -\frac{2l}{N_t} -A}{B}\right]\right)
\end{align}
%
which after some algebra simplifies to
%
\begin{align}\label{eq:last_step_gnm_orthonorm}
\sum_{nm}&\tilde{g}_{nm}[l]\tilde{g}_{nm}^*[l']=\delta_{ll'}\left(1+\cos\left(\frac{\pi}{2}\frac{1-2A}{B}\right)\cos\left(2\frac{\pi}{2}\frac{\frac{2l}{N_t}-\bar m -A}{B}-\frac{\pi}{2}\frac{1-2A}{B}\right)\right)=\delta_{ll'}
\end{align}
%
where we used the fact of $2A+B=1$. Therefore we find that for every value of $l$ we have%
%
\begin{align}
\sum_{nm}&\tilde{g}_{nm}[l]\tilde{g}_{nm}^*[l']=\delta_{ll'}
\end{align}
%
\subsection{The second (quasi)normality condition}
We want to evaluate
%
\begin{align}
\sum_{l}&\tilde{g}_{nm}[l]\tilde{g}_{pq}^*[l]
\end{align}
%
we immediatly note that, since the window support in frequency is $N_t/2$ we are forced to have $q=m$. Now we have to consider three cases.

\textbf{Case $m=0$}, for which we have
%
\begin{align}
\sum_{l}\tilde{g}_{n0}[l]\tilde{g}_{p0}^*[l]&=\frac{1}{2}\sum_{l=-N/2}^{N/2-1} e^{-4\pi i l (n-p)/N_t}\, \tilde{\varphi}^2[l]=\frac{1}{2}\sum_{l=-N_t/2}^{N_t/2-1} e^{-4\pi i l (n-p)/N_t}\, \tilde{\varphi}^2[l]=\nonumber\\
&=\sum_{l=-N_t/2}^{-1}\left(e^{-4\pi i l (n-p)/N_t}\tilde{\varphi}^2[l]+e^{-4\pi i (l+N_t/2) (n-p)/N_t}\tilde{\varphi}^2[l+N_t/2]\right)
\end{align}
%
but now we note that $e^{-4\pi i (l+N_t/2) (n-p)/N_t}=e^{-4\pi i l(n-p)/N_t}$ and in the previous subsection of this appendix we have also shown that for $|l|\leq N_t/2$ we have $\tilde{\varphi}^2[l]+\tilde{\varphi}^2[l+N_t/2]=\frac{2}{N_t}$, therefore the last expression simplifies to
%
\begin{align}
\sum_{l}\tilde{g}_{n0}[l]\tilde{g}_{p0}^*[l]&=\frac{1}{N_t}\sum_{l=-N_t/2}^{-1}e^{-4\pi i l (n-p)/N_t}
\end{align}
%
this is a geometric series in $l$, which vanishes for every choice of $n-m$, but the only two cases $n-m=0$ and $n-m=\pm N_t/2$, for which the exponential evaluates to $1$, therefore obtaining, for the first case
%
\begin{align}
\sum_{l}\tilde{g}_{n0}[l]\tilde{g}_{p0}^*[l]&=\frac{1}{2}\left(\delta_{n,p}+\delta_{n,p\pm \frac{N_t}{2}}\right)
\end{align}

\textbf{Case $m=N_f$}, for which we have
%
\begin{align}
\sum_{l}\tilde{g}_{n N_f}[l]\tilde{g}_{p N_f}^*[l]&=\frac{1}{2}\sum_{l=-N/2}^{N/2-1} e^{-4\pi i l (n-p)/N_t}\, \left(\tilde{\varphi}[l-N/2]+\tilde{\varphi}[l+N/2]\right)^2=\nonumber\\
&=\frac{1}{2}\sum_{l=-N/2}^{N/2-1} e^{-4\pi i l (n-p)/N_t}\, \left(\tilde{\varphi}^2[l-N/2]+\tilde{\varphi}^2[l+N/2]\right)=\nonumber\\
&=\frac{1}{2}\sum_{l=-N}^{-1} e^{-4\pi i (l+N/2)(n-p)/N_t}\tilde{\varphi}^2[l]+\frac{1}{2}\sum_{l=0}^{N-1} e^{-4\pi i (l-N/2)(n-p)/N_t}\tilde{\varphi}^2[l]=\nonumber\\
&=\frac{1}{2}\sum_{l=-N_t/2}^{-1} e^{-4\pi i l(n-p)/N_t}\tilde{\varphi}^2[l]+\frac{1}{2}\sum_{l=0}^{N_t/2-1} e^{-4\pi i l(n-p)/N_t}\tilde{\varphi}^2[l]=\nonumber\\
&=\frac{1}{2}\sum_{l=-N_t/2}^{N_t/2-1} e^{-4\pi i l(n-p)/N_t}\tilde{\varphi}^2[l]
\end{align}
%
and from here the proof is identical to the $m=0$ case, leading to
%
\begin{align}
\sum_{l}\tilde{g}_{nN_f}[l]\tilde{g}_{pN_f}^*[l]&=\frac{1}{2}\left(\delta_{n,p}+\delta_{n,p\pm \frac{N_t}{2}}\right)
\end{align}

\textbf{Case $0<m<N_f$}, for which we have
%
\begin{align}
\sum_{l}\tilde{g}_{nm}[l]\tilde{g}_{pm}^*[l]&=\frac{1}{2}\sum_{l=-N/2}^{N/2-1} \Bigg[e^{-2\pi i l (n-p)/N_t}\,\Big(C_{nm}\,\tilde{\varphi}[l - mN_t/2] + C_{nm}^*\,\tilde{\varphi}[l + mN_t/2]\Big) \\ 
& \qquad \times \Big(C_{pm}^*\,\tilde{\varphi}[l - mN_t/2] + C_{pm}\,\tilde{\varphi}[l + mN_t/2]\Big) \Bigg]\nonumber\\
&=\frac{1}{2}\sum_{l=-N/2}^{N/2-1} e^{-2\pi i l (n-p)/N_t}\,\Big(C_{nm}C^*_{pm}\,\tilde{\varphi}^2[l - mN_t/2] + C_{nm}^*C_{pm}\,\tilde{\varphi}^2[l + mN_t/2]\Big)
\end{align}
%
Let's take the first term and shift the frequency, rewriting
%
\begin{align}
\frac{1}{2}C_{nm}C^*_{pm}\sum_{l=-N_t/2}^{N_t/2-1} e^{-2\pi i (l+mN_t/2) (n-p)/N_t}\tilde{\varphi}^2[l]=\frac{1}{2}C_{nm}C^*_{pm}e^{-\pi i m(n-p)}\sum_{l=-N_t/2}^{N_t/2-1} e^{-2\pi i l(n-p)/N_t}\tilde{\varphi}^2[l]
\end{align}
%
and also in this case, from here the proof is identical to the $m=0$ case, but where in this case, given the different exponent, the series is nonvanishing for $n=p$ and not for $n=p\pm N_t/2$. The expression then simplifies to
%
\begin{align}
\frac{1}{2}C_{nm}C^*_{pm}\sum_{l=-N_t/2}^{N_t/2-1} e^{-2\pi i (l+mN_t/2) (n-p)/N_t}\tilde{\varphi}^2[l]=\frac{1}{2}C_{nm}C^*_{nm}=\frac{1}{2}
\end{align}
%
an analogous proof holds for the term with $+mN_t/2$, therefore the full result is that in this case, for $0<m<N_f$, we find
%
\begin{align}
\sum_{l}\tilde{g}_{nm}[l]\tilde{g}_{pm}^*[l]=\delta_{n,p}
\end{align}
%
To wrap up all the cases we find
%
\begin{align}
\sum_{l}\tilde{g}_{nm}[l]\tilde{g}_{pq}^*[l]=\delta_{m,q}\begin{cases}
    \frac{1}{2}\left(\delta_{n,p}+\delta_{n,p\pm \frac{N_t}{2}}\right) \qquad m=\{0,N_f\}\\
    \delta_{n,p} \qquad 0<m<N_f
\end{cases}
\end{align}




\subsection{Normalization of the Meyer window}\label{app:normalization_Meyer}

We verify that the Meyer window defined in
Eq.~\eqref{eq:phi_discrete} (the discrete case of \eqref{meyer_window_discrete}) is correctly normalized. We compute
%
\begin{align}
\sum_{l=-N/2}^{N/2-1} \tilde{\varphi}^2[l]&=2\sum_{l=0}^{A\frac{N_t}{2}-1}\frac{2}{N_t}-\frac{2}{N_t}+2\sum_{l=A\frac{N_t}{2}}^{(A+B)\frac{N_t}{2}-1}\frac{2}{N_t}\cos^2\left[\frac{\pi}{2}\frac{\frac{2l}{N_t}-A}{B}\right]=2A-\frac{2}{N_t}+\frac{4}{N_t}\sum_{k=0}^{B\frac{N_t}{2}-1}\cos^2\left[\frac{\pi}{2}\frac{2k}{BN_t}\right]=\nonumber\\
&=2A-\frac{2}{N_t}+\frac{4}{N_t}\sum_{k=0}^{B\frac{N_t}{2}-1}\frac{1}{2}+\frac{4}{N_t}\sum_{k=0}^{B\frac{N_t}{2}-1}\frac{1}{2}\cos\left(\frac{2\pi k}{BN_t}\right)=2A-\frac{2}{N_t}+B+\frac{2}{N_t}=1
\end{align}
%
where on the first line the term $-\frac{2}{N_t}$ avoids double counting $l=0$, where the last summation equals $+\frac{2}{N_t}$ after performing the geometrical series and where we remember that $2A+B=1$ is a property of the window function.

\section{Detailed derivation for the forward and inverse transforms}\label{app:forward_derivation}

% \ob{Again, focusing on $d = 1$. Can we not do this for $d > 1$ ? I feel like it should be possible. Reading up on the incomplete beta function, we have the following identity : $\nu_{d}(1 - y) + \nu_{d}(y) = 1$. This means that we can simplify expressions like:
% \begin{align}
%     \cos^2\left(\frac{\pi}{2}\nu_d(y)\right) + \cos^2\left(\frac{\pi}{2}\nu_d(1 - y)\right) 
%     =  \cos^{2}\left(\frac{\pi}{2}\nu_d(y)\right) + \cos^2\left(\frac{\pi}{2}\left[1 - \nu_d(y)\right]\right) = \cos^{2}\left(\frac{\pi}{2}\nu_d(y)\right) + \sin^2\left(\frac{\pi}{2} \nu_d(y)\right) = 1
% \end{align}
% Here we have used $\cos(\pi/2 - y) = \sin(y)$ and $\nu_d(1 - y) = 1 - \nu_d(y)$. I feel like this identity will help with some of the computations above?
% }
% \ob{OK! let me take a look. }
% \gm{If you remember I derived the window normality for generic $d$ some time ago, using basically your idea. But I did not want to introduce complications that may distract from the main message of these calculations.}
 This appendix provides the full algebraic steps leading to the
compact expressions for the forward WDM coefficients presented in
Section~\ref{sec:wdm}. We work in the Fourier domain throughout,
writing the forward transform as
\begin{align}
  w_{nm}=\sum_{l=-N/2}^{N/2-1} \tilde{x}_l\,\tilde{g}_{nm}^*(f[l])\,,
\end{align}
where the discrete frequencies are mapped as $f[l]=l\,\Delta f=l/T$.

\subsection{Interior channels ($0<m<N_f$)}

Inserting the interior-channel basis element from
Eq.~\eqref{eq:gnm_full_discrete} and splitting the sum over positive and
negative frequencies gives
\begin{align}
w_{nm}&=\frac{1}{\sqrt{2}}\sum_{l=0}^{N/2-1} \tilde{x}_l\, e^{2\pi iln/N_t}\left[C_{nm}^*\tilde{\varphi}\!\left(\!(l-\tfrac{mN_t}{2})\Delta f\right)+C_{nm}\tilde{\varphi}\!\left(\!(l+\tfrac{mN_t}{2})\Delta f\right)\right]\nonumber\\
&+\frac{1}{\sqrt{2}}\sum_{l=-N/2}^{-1} \tilde{x}_l\, e^{2\pi iln/N_t}\left[C_{nm}^*\tilde{\varphi}\!\left(\!(l\!-\!\tfrac{mN_t}{2})\Delta f\right)+C_{nm}\tilde{\varphi}\!\left(\!(l\!+\!\tfrac{mN_t}{2})\Delta f\right)\right].
\end{align}
Using the reality
condition $\tilde{x}_{-l} = \tilde{x}_l^*$ together with the
symmetry $\tilde{\varphi}(-f) = \tilde{\varphi}(f)$, the second sum
is the complex conjugate of the first. The boundary terms at $l=0$
and $l=-N/2$ vanish because the window support does not extend to
those frequencies for interior channels. We therefore obtain
\begin{align}\label{eq:app_wnm_general}
  w_{nm}&=\sqrt{2}\,\Re\sum_{l=1}^{N/2-1} \tilde{x}_l\, e^{2\pi
  iln/N_t}\,C_{nm}^*\,\tilde{\varphi}\!\left(\!(l-\tfrac{mN_t}{2})\Delta f\right)\nonumber\\
  &=\sqrt{2}\,(-1)^{nm}\,\Re\sum_{l'=-N_t/2}^{N_t/2}
  \tilde{x}_{l'+mN_t/2}\, e^{2\pi
  il'n/N_t}\,C_{nm}^*\,\tilde{\varphi}(l'\,\Delta f)\,,
\end{align}
where the last equality follows from the substitution $l' = l -
mN_t/2$ and the fact that $\tilde{\varphi}(l'\,\Delta f)$ is nonzero only
for $|l'| \leq N_t/2$. This is Eq.~\eqref{eq:wnm_interior} of the
main text, and corresponds (up to normalization conventions) to
Eq.~17 of Cornish~\cite{Cornish:2020odn}.

\subsection{DC channel ($m = 0$)}

For the DC edge channel, the basis element has a doubled time-shift
exponent and no $C_{nm}$ factor. The same positive/negative frequency
splitting yields
\begin{align}
  w_{n0}&=\frac{1}{\sqrt{2}}\sum_{l=0}^{N/2-1} \tilde{x}_l\, e^{4\pi
  iln/N_t}\,\tilde{\varphi}(l\,\Delta f)
  +\frac{1}{\sqrt{2}}\sum_{l=-N/2}^{-1} \tilde{x}_l\, e^{4\pi
  iln/N_t}\,\tilde{\varphi}(l\,\Delta f)
  \nonumber\\
  &=\,\sqrt{2}\Re\left[\sum_{l=1}^{N_t/2} \tilde{x}_{l}\, e^{4\pi
  iln/N_t}\,\tilde{\varphi}(l\,\Delta f)\right]+\frac{1}{\sqrt{2}}\tilde{x}_0\,\tilde{\varphi}(0)\,,
\end{align}
which is Eq.~\eqref{eq:wn0} of the main text.
\subsection{Nyquist channel ($m=N_f$)}

\begin{align}
w_{nm}&=\frac{1}{\sqrt{2}}\sum_{l=0}^{N/2-1} \tilde{x}_l\, e^{2\pi iln/N_t}\left[\tilde{\varphi}\!\left(\!(l-\tfrac{N}{2})\Delta f\right)+\tilde{\varphi}\!\left(\!(l+\tfrac{N}{2})\Delta f\right)\right]+
\frac{1}{\sqrt{2}}\sum_{l=-N/2}^{-1} \tilde{x}_l\, e^{2\pi iln/N_t}\left[\tilde{\varphi}\!\left(\!(l\!-\!\tfrac{N}{2})\Delta f\right)+\tilde{\varphi}\!\left(\!(l\!+\!\tfrac{N}{2})\Delta f\right)\right]\nonumber\\
&=\frac{1}{\sqrt{2}}\sum_{l=0}^{N/2-1} \tilde{x}_l\, e^{2\pi iln/N_t}\left[\tilde{\varphi}\!\left(\!(l-\tfrac{N}{2})\Delta f\right)+\tilde{\varphi}\!\left(\!(l+\tfrac{N}{2})\Delta f\right)\right]+
\frac{1}{\sqrt{2}}\sum_{l=1}^{N/2} \tilde{x}_{-l}\, e^{-2\pi iln/N_t}\left[\tilde{\varphi}\!\left(\!(-l\!-\!\tfrac{N}{2})\Delta f\right)+\tilde{\varphi}\!\left(\!(-l\!+\!\tfrac{N}{2})\Delta f\right)\right].
\end{align}
%
the $l=0$ term in the first summation vanishes, and the $l=N/2$ term is treated separately:
%
\begin{align}
  w_{nm}&=\sqrt{2}\,\Re\sum_{l=1}^{N/2-1} \tilde{x}_l\, e^{2\pi
  iln/N_t}\,\,\tilde{\varphi}\!\left(\!(l-\tfrac{N}{2})\Delta f\right)+\frac{1}{\sqrt{2}}\tilde{x}_{N/2}\, e^{-\pi
  inN_f}\tilde{\varphi}(0)\nonumber\\
  &=\sqrt{2}\,\,\Re\sum_{l'=-N_t/2}^{-1}
  \tilde{x}_{l'+N/2}\, e^{2\pi
  il'n/N_t}\,\tilde{\varphi}(l'\,\Delta f)+\frac{1}{\sqrt{2}}\tilde{x}_{N/2}\, \tilde{\varphi}(0)\,,
\end{align}
%
assuming an even value for $N_f$.

\subsection{The inverse transform}\label{app:inverse_transform}

We write the inverse transform of the $w_{nm}$ coefficients to recover the signal $x[l]$ for $l\geq 0$ since we know that for $l<0$ the reality condition holds ($x[-l]=x^*[l]$):
%
\begin{align}
x[l]=\sum_{n=0}^{N_t-1}\sum_{m=0}^{N_f}w_{nm}g_{nm}[l]&=\frac{1}{\sqrt{2}}\sum_{n}\sum_{0<m<N_f}w_{nm}e^{-2\pi iln/N_t}\left[C_{nm}\tilde{\varphi}\!\left(\!(l-\tfrac{mN_t}{2})\Delta f\right)+C^*_{nm}\tilde{\varphi}\!\left(\!(l+\tfrac{mN_t}{2})\Delta f\right)\right]\nonumber\\
&+\frac{1}{\sqrt{2}}\sum_{n}w_{n0}e^{-4\pi iln/N_t}\tilde{\varphi}(l\,\Delta f)+\frac{1}{\sqrt{2}}\sum_{n}w_{nN_f}e^{-4\pi iln/N_t}\left[\tilde{\varphi}\!\left(\!(l-\tfrac{N}{2})\Delta f\right)+\tilde{\varphi}\!\left(\!(l+\tfrac{N}{2})\Delta f\right)\right]=\nonumber\\
&=\begin{cases}
x_0[l]\qquad\qquad\qquad \qquad l<N_t/2\\
x_m[l]\qquad\qquad\qquad m N_t/2\leq l<(m+1)N_t/2\\
x_{N_f}[l]\qquad\qquad\qquad l \geq N/2-N_t/2
\end{cases}
\end{align}
%
where we note that for each value of $l$, the inverse transform depends on which of the $m$ bands the $l$ value belongs, in particular
%
\begin{align}
x_0[l]&=\frac{1}{\sqrt{2}}\sum_{n}w_{n1}e^{-2\pi iln/N_t}C_{n1}\tilde{\varphi}\!\left(\!(l-\tfrac{N_t}{2})\Delta f\right)+\frac{1}{\sqrt{2}}\sum_{n}w_{n0}e^{-4\pi iln/N_t}\tilde{\varphi}(l\,\Delta f)\nonumber\\
x_m[l]&=\frac{1}{\sqrt{2}}\sum_n w_{nm}e^{-2\pi iln/N_t}C_{nm}\tilde{\varphi}\!\left(\!(l-\tfrac{mN_t}{2})\Delta f\right)+\frac{1}{\sqrt{2}}\sum_n w_{n(m+1)}e^{-2\pi iln/N_t}C_{n(m+1)}\tilde{\varphi}\!\left(\!(l-\tfrac{(m+1)N_t}{2})\Delta f\right)\nonumber\\
x_{N_f}[l]&=\frac{1}{\sqrt{2}}\sum_n w_{n(N_f-1)}e^{-2\pi iln/N_t}C_{n(N_f-1)}\tilde{\varphi}\!\left(\!(l-\tfrac{N}{2}+\tfrac{N_t}{2})\Delta f\right)+\frac{1}{\sqrt{2}}\sum_n w_{n N_f}e^{-4\pi iln/N_t}\tilde{\varphi}\!\left(\!(l-\tfrac{N}{2})\Delta f\right)
\end{align}
Which can be written in a slightly more convenient form:

\begin{align}
\tilde{x}_m[l] = \dfrac{1}{\sqrt{2}}
\begin{cases}
\tilde{\varphi}[l]\, \displaystyle\sum_{n=0}^{N_t - 1} w_{n0}\, e^{-4\pi i n l / N_t},
& m = 0,\ \ 0 \le l < N_t/2, \\[10pt]
\tilde{\varphi}[l - mN_t/2]\, \displaystyle\sum_{n=0}^{N_t - 1} C_{nm}\, w_{nm}\, e^{-2\pi i n l / N_t},
& 1 \le m \le N_f - 1,\ \ (m{-}1)\tfrac{N_t}{2} \le l < (m{+}1)\tfrac{N_t}{2}, \\[10pt]
\tilde{\varphi}[l - N/2]\, \displaystyle\sum_{n=0}^{N_t - 1} w_{nN_f}\, e^{-4\pi i n l / N_t},
& m = N_f,\ \ N/2 - N_t/2 \le l \le N/2,
\end{cases}
\end{align}
which is presented in the main text \eqref{eq:inverse_transform_discrete}. 

\section{WDM likelihood and the Local Stationary Assumption}\label{app:stats}
\subsection{The Evolutionary PSD}

The Frequency domain is an excellent domain of choice if the underlying noise process is stationary with the extra feature of being circulant. Such conditions allows the time-domain auto-covariance matrix to be diagonalised in a Fourier basis facilitating a rapid to evaluate likelihood function, essential for parameter inference. If the noise were non-stationary, the frequency domain noise components exhibit correlations making (1) the noise covariance matrix expensive to invert and, potentially, non-singular if apodization methods are applied.

Small time-dependent drifts to the PSD are problematic for frequency domain analyses. One frequency domain approach is segmenting the data on time-scales similar to the non-stationary effects and performing inference on each chunk then assuming independence, treating the log-likelihood as a sum of smaller segment log-likelihoods. In time-frequency, we can show that even with slight non-stationary features that assuming a ``diagonal" based WDM noise covariance matrix should still be sufficient for unbiased parameter inference. 

Let us define the discrete time-domain auto-covariance function for a stochastic (real) process $x$ via
\begin{equation}\label{eq:general_autocov_time}
R[k,k'] = \mathbb{E}[n[k],n[k']]
\end{equation}

for time-bins $k,k'\in\{0,\ldots,N-1\}$. For the general auto-covariance function \eqref{eq:general_autocov_time}, the stationary assertion means that 
\begin{equation}
    R[k,k'] = R[0,k - k'] = \mathbb{E}[n[0]n[k - k']] = R(\kappa)
\end{equation}
for $\kappa = k - k'$ a lag parameter. Such a condition forces the time-domain noise covariance matrix to obey a Toeplitz like structure (not necessarily circulant). The Wiener-Khinchin theorem states that the Fourier transform of the time-domain auto-covariance function is related to the processes power spectral density

\begin{equation}
S[l] \equiv 2\Delta t\sum_{\kappa = 0}^{N - 1}R(\kappa)e^{-2\pi i l \kappa/N}\,.
\end{equation}
The factor of 2 results from using a one-sided PSD and $\Delta t$ from the sum. It is then possible to build the frequency domain noise covariance matrix via 
\begin{align}
(\tilde{\boldsymbol{\Sigma}})_{l,l'}&=\mathbb{E}[\tilde{x}[l],\tilde{x}^\star[l']]  = \sum_{k,k'}\mathbb{E}[x[k] x[k']]e^{-2\pi i (kl - k'l')/N}\\
&= \frac{1}{2\Delta t}\sum_{\bar{k}}2\Delta t\sum_{\kappa}R[\bar{k} + \kappa/2, \bar{k} - \kappa/2]e^{-2\pi i \left[\frac{l + l'}{2} \right]\kappa/N}e^{-2\pi i \bar{k}(l - l')/N} \\
&= \frac{1}{2\Delta t} \sum_{\bar{k}} S\left[\bar{k}, \frac{1}{2}(l + l')\right]e^{-2\pi i \bar{k}(l - l')/N}\,,
\end{align}
%
where $\bar k=\frac{k+k'}{2}$.
This last formula is also known as Wigner-Ville representation \cite{PhysRev.40.749}. Notice for arbitrary values of $\kappa$ the frequency domain noise covariance matrix now has non-trivial off diagonal elements that must be accounted for during inference. Notice that if we drop $\bar{k}$ from the evolutionary PSD and evaluating the sum over the complex exponential, we obtain the familiar stationary expression
\begin{equation}\label{app:eq_stat_covariance_normalisation}
(\tilde{\boldsymbol{\Sigma}})_{l,l'} = \frac{NS[l]}{2\Delta t}\delta_{ll'}\,.
\end{equation}
%
\ob{Careful Ollie -- Toeplitz + Circulant. I think what we need to say here is that $R$ encodes \emph{all} of the statistical information about the process $x$, therefore being circulant (periodic statistics at end points).}

\subsection{The WDM Covariance Matrix}\label{app:WDM_cov_calcs}

Consider now the expression for the WDM covariance matrix
\begin{align}
(\boldsymbol{\Sigma}_w)_{(n,m),(p,q)} &= \sum_{l,l'}\tilde{g}_{nm}^{\star}[l]\tilde{g}_{pq}[l']\mathbb{E}[\tilde{x}^\star[l]\tilde{x}[l']] \\
& = \frac{1}{2\Delta t}\sum_{l,l'}\tilde{g}_{nm}^{\star}[l]\tilde{g}_{pq}[l'] \sum_{\bar{k}} S\left[\bar{k}, \frac{1}{2}(l + l')\right]e^{-2\pi i \bar{k}(l - l')/N}\,, \label{eq:WDM_cov_general_app}
\end{align}
%
We will now look at two separate cases -- the leading diagonal $(n,m) = (p,q)$ and off diagonal elements $(n,m) \neq (p,q)$. The atoms from the Wilson basis $\tilde{g}$ are peaked around the value of $l \approx mN_t/2$. This is where the support of the Wilson basis lies. This means that the dominant contribution of the sum are for those $l,l'$ in the vicinity of $mN_t/2$ and $qN_t/2$. Let us now assume that the evolutionary PSD varies slowly over the regime $(l + l')/2 \approx mN_t/2$ and the temporal time index $\bar{k} \approx nN_f$

Therefore, let us perform a discrete expansion for both the time and frequency indices in $S[\bar{k}, (l + l')/2]$ in this vicinity and applying the \emph{local stationary approximation}
%
\begin{equation}\label{eq:local_stat_approx}
S\left[\bar{k}, \frac{1}{2}(l + l')\right] \approx S[nN_f, mN_t/2] + \left(\frac{l + l'}{2} - mN_t/2\right)\partial_{f}S + (\bar{k} - nN_f)\partial_{t}S + \text{corrections}
\end{equation}
The expansion in \eqref{eq:local_stat_approx} is only valid under the two conditions
\begin{equation}
\Delta T \big |\partial_{t}\log S(t,f)\big |_{(t,f) = (t_0,f_0)} \ll 1 \,, \quad \text{and} \quad \Delta F \big |\partial_{f}\log S(t,f)\big |_{(t,f) = (t_0,f_0)} \ll 1\,. 
\end{equation}
The first condition is the true local stationary approximation, where we assume that the logarithmic sensitivity of the PSD $\partial_{t}\log S(t,f)$ varies slowly over a wavelet bandwidth $\Delta T$. Such a condition is violated if the process exhibits major $t$ dependent drifts in the noise statistics. The second condition implies that the noise process over a wavelet frequency band $\Delta F$ smoothly evolves as a function of logarithmic frequency $\partial_{f}\log S(t,f)$. Minima in the evolutionary PSD as a function of frequency enforced via Fourier transforms of delays (turning points) as a result of Time Delay Interferometric computations could violate such an assumption. The local stationary approximation is an approximation that should therefore not be taken lightly and the analyst should be careful when envoking it \ob{improve}. 
\subsubsection{Leading Diagonal}
Now, for the leading diagonal, substituting the leading order term from \eqref{eq:local_stat_approx} into \eqref{eq:WDM_cov_general_app} and evaluating the sum, one obtains
\begin{align}
    (\boldsymbol{\Sigma}_w)_{(n,m),(p,q)}
& \approx \frac{1}{2\Delta t} S[nN_f,mN_t/2]\sum_{l,l',\bar{k}}\tilde{g}_{nm}^{\star}[l]\tilde{g}_{nm}[l'] e^{-2\pi i \bar{k}(l - l')/N}\,, \\
& =  \frac{N}{2\Delta t} S[nN_f,mN_t/2] \sum_{l}|\tilde{g}_{nm}[l]|^2 \\
& = \frac{N}{2\Delta t} S[n,mN_t/2]\,.
\end{align}
Now let's look at teach part of the first order contributions in the evolutionary PSD
\begin{align}
&\frac{1}{2\Delta t}\partial_{f}S[nN_f,m N_t/2]\sum_{l,l',\bar{k}}\tilde{g}_{nm}^{\star}[l]\tilde{g}_{nm}[l'] \left(\frac{l + l'}{2} - mN_t/2\right) e^{-2\pi i \bar{k}(l - l')/N)}\,, \\
&= \frac{N}{2\Delta t}\partial_{f}S[nN_f,m N_t/2]\sum_{l}|\tilde{g}_{nm}[l]|^{2}\left[l - mN_t/2\right] \,, \\
& = 0?
\end{align}
\ob{I think this is zero because of the Meyer window properties. Can you check Giorgio? }
\gm{Yes, for $m>0$ you have that 
\begin{align}
|\tilde g_{nm}[l]|^2=\frac{1}{2}\left(\tilde{\varphi}^2[l - mN_t/2]+\tilde{\varphi}^2[l + mN_t/2]\right)
\end{align}
and now if you do the substitution $l\to l-mN_t/2$ for the first term, and $l\to l+mN_t/2$ for the second, you realize that they are both integrals of odd functions on a symmetric domain. Similar proof for $m=0,N_f$. Again, we probably should be careful in understanding under which conditions we can Taylor-expand in frequency...
}\ob{Magic, can you argue this below for me? }
\gm{Basically, the Taylor works if the frequency spectrum does not vary much through the frequency bins, i.e. $\frac{dS}{df}\frac{\Delta F}{S}\ll1$, right?}
The support of $\tilde{g}$ is in the vicinity of $l = mN_t/2$. Loosely speaking, when $\tilde{g}$ is maximized, it is then suppressed by the factor $l - mN_t/2$. The temporal term takes the form:
\begin{align}
&\frac{1}{2\Delta t}\partial_{t}S[nN_f,m N_t/2]\sum_{l,l',\bar{k}}\tilde{g}_{nm}^{\star}[l]\tilde{g}_{nm}[l'] \left(\bar{k} - nN_f\right) e^{-2\pi i \bar{k}(l - l')/N}\,, \\
&= \frac{1}{2\Delta t}\partial_{t}S[nN_f,m N_t/2]\sum_{\bar{k}}\left[\left(\sum_{l}\tilde{g}_{nm}^{\star}[l]e^{2\pi i \bar{k}l/N}\sum_{l'}\tilde{g}_{nm}[l']e^{-2\pi i l' \bar{k} / N}\right)\left(\bar{k} - nN_f\right)\right] \\
& = \frac{N^2}{2\Delta t}\partial_{t}S[nN_f,m N_t/2]\sum_{\bar{k}}|g_{nm}[\bar{k}]|^2 (\bar{k} - nN_f) \,,
\end{align}
Where we have identified inverse Fourier transforms of the $\tilde{g}$ basis functions to dramatically simplify the expression above. 

From Necula \cite{Necula_2012}, the time-domain Wilson basis follows from inverse-DFT of Eq.~(\ref{eq:gnm_full_discrete}):
\begin{equation}
    g_{nm}[\bar{k}] = \sqrt{2}\,\varphi[\bar{k} - nN_f]\,\cos\!\left(\frac{\pi m \bar{k}}{N_f} + \theta_{nm}\right),
    \label{eq:wilson_time_domain}
\end{equation}
\gm{I am not sure this formula is ok, I find a bit differently but maybe I'm wrong. The important thing is that when you do the $||^2$ you get an odd function integrated on a symmetric time domain. For $d=1$ the explicit $\varphi$ expression is the last equation of \url{https://en.wikipedia.org/wiki/Meyer_wavelet}, but we should be pedantic and do the proper calculation, as in the frequency domain. Also in this case, I believe we have to ask $\frac{dS}{dt}\frac{\Delta T}{S}\ll1$ to have the Taylor expansion valid.}\ob{I had a sign wrong... did that fix things?}
for the alternating phase factor $\theta_{nm}$
\begin{equation}
\theta_{nm} = \begin{cases} 0 & (n + m) \ \text{even} \\
\pi/2 & (n + m) \ \text{odd} 
\end{cases}
\end{equation}
where $\varphi$ is the inverse DFT of the Meyer window $\tilde\varphi$. The atom is localised in time at the centre $k = nN_f$ with width set by the inverse-DFT of $\tilde\varphi$, and modulated at the channel frequency $m\Delta F$. 

In a similar manner to the frequency domain computation above, the quantity
\begin{equation}
\sum_{\bar{k}}|g_{nm}[\bar{k}]|^2 (\bar{k} - nN_f) \approx \frac{1}{\sqrt{2}}\sum_{\bar{k}}(\bar{k} - nN_f)|\varphi[\bar{k}]|^2 \approx 0
\end{equation}
since $\varphi$ is an even function $\varphi[-k] = \varphi[k]$, around it's centre $nNf$. This implies that the temporal piece of the evolutionary power spectral density also vanishes. The diagonal components therefore take the form 
\begin{equation}
(\boldsymbol{\Sigma}_w)_{(n,m),(n,m)} = \frac{N}{2\Delta t}S[nN_f, mN_t/2]\left[1 + \mathcal{O}(\text{quadratic corrections})\right]
\end{equation}
in other words, in the locally stationary limit. 

\subsubsection{off diagonals}


We leave the off-diagonal elements to another day. 




In the very simple case of wide sense stationary noise, $\partial_{t}S = 0$, the pre-factor (involving sums) of $\partial_{f}S$ is still zero. It can then be shown that
\begin{align}
    (\boldsymbol{\Sigma}_w)_{(n,m),(p,q)}
& \approx \frac{1}{2\Delta t} S[nN_f,mN_t/2]\sum_{l,l',\bar{k}}\tilde{g}_{nm}^{\star}[l]\tilde{g}_{pq}[l'] e^{-2\pi i \bar{k}(l - l')/N}\,, \\
& =  \frac{N}{2\Delta t} S[nN_f,mN_t/2] \sum_{l}\tilde{g}_{nm}^{\star}[l]\tilde{g}_{pq}[l']\\
& = \frac{N}{2\Delta t} S[nN_f,mN_t/2]\delta_{np}\delta_{mq}\\
& = \delta_{mq} S[nN_f, mN_t/2] \cdot \begin{cases}  \delta_{np} & m \in \{1,\ldots, N_f - 1\} \\
\frac{1}{2}\left(\delta_{n,p} + \delta_{n, p \pm N_t/2}\right) & m \in \{0, N_f\}\,.
\end{cases}
\end{align}
as required.

the evolutionary PSD $S^{\rm WDM}_{nm}$-- the noise PSD with the freedom to evolve as a function of both frequency  $f_m$ and time $t_n$. We remark that the evolutionary PSD is \emph{not} a covariance matrix, but only when combined with the Kronecker deltas above as a result of local-stationarity can it be used as a covariance matrix. The factor $N N_f$ relative to the natural pixel variance $S_n \Delta F$ arises from our discrete-DFT convention (Eq.~\ref{eq:dft-definition}), in which $\tilde{x}[l]$ does not carry an explicit factor of $\Delta t$. Cornish~\cite{Cornish:2020odn} uses $\Delta t$-normalised transforms and accordingly obtains $S^{\mathrm{WDM}}_{nm} \approx S_n(f_m)\,\Delta F$.

If we include the zeroth and Nyquist frequency bins $ m = \{0,N_f\}$ and using the first equation in \eqref{eq:orthonormality}, we see that the WDM covariance at these bins takes a slightly different form
\begin{equation}\label{eq:cov_zeroth_nyquist_WDM}
(\boldsymbol{\Sigma}_w)^{m \in \{0, N_f\}}_{\alpha\beta} = \frac{1}{2}\, S^\text{WDM}_{nm}\, \delta_{mq}\bigl(\delta_{np} + \delta_{n,\,p\pm N_t/2}\bigr).
\end{equation}
%
implying that the variance of the zeroth and Nyquist bins is different from the interior bins. This is analogous to the frequency domain analogue as described in Appendix A of \cite{Burke:2025bun}. 

% One can write the likelihood cleanly over all frequency bins only if the evolutionary spectrum $S^{\rm WDM}_{nm}$ is assumed known. In this case, the likelihood is proportional  to the quadratic piece including only the signal parameters of interest. However, for noise estimate, the logarithmic term must be included to account for variations in evolutionary PSD. In short -- if one wishes to be pedantic and write over all frequency bins then the likelihood is much more complicated and requires the use of the pseudo determinant to evaluate the $\log \det$ term -- for very little gain. 

Over the interior bins only, substituting the diagonal covariance~\eqref{eq:cov_interior_WDM} into
the general Gaussian likelihood yields the WDM Whittle-based log-likelihood:
\begin{align}\label{eq:wdm_whittle}
  \ln p(d|h) &\propto -\sum_{\substack{n \ \\ m\ \in\rm \ int}}\bigg[\ln\left(\frac{1}{N_tN_f} S^{\text{WDM}}_{nm}\right) \\
  & \qquad \qquad \qquad \qquad \quad +\frac{1}{2}  
  \frac{(d_{nm} - h_{nm})^2}{S^{\text{WDM}}_{nm}}\bigg],
\end{align}
\gm{Is there a - sign missing in front of $\frac{1}{2}$ in this equation above? Also, why did you split $N$ to $N_t N_f$?}\ob{So everything is in terms of WDM constants $N_t,N_f$. It's not necessary, just preferred (by me).}
\gm{Also $N$ technically is a WDM constant...but let's decide the others}
where $d_{nm}$, $h_{nm}$ are the WDM coefficients of the data and
signal model, respectively. The likelihood \eqref{eq:wdm_whittle} is sufficient for both noise and signal parameter estimation. From this point onwards, we assume a known noise covariance matrix in the WDM domain, implying that the likelihood \eqref{eq:wdm_whittle} is proportional up to the quadratic piece (second term). 

In the ideal orthogonal limit, whitened noise --
noise with $S_n[l] \propto $ const. across the support of each
atom -- reduces this covariance to diagonal form. In that case the WDM
coefficients are uncorrelated, and for Gaussian noise they are
therefore independent. In the sampled packed implementation used here,
it is useful to distinguish two statements. First, the interior atoms
($m=1,\ldots,N_f-1$) satisfy the discrete orthonormality relation to
machine precision. Second, the packed real-valued storage includes
special DC and Nyquist edge channels, so the full stored coefficient
array is best viewed as a convenient packed representation rather than
a minimal orthonormal basis.
For this reason, the decorrelation discussion below focuses on the
interior coefficients. For whitened noise their empirical
correlations remain very small, while for coloured, non-stationary, or
non-Gaussian noise off-diagonal structure becomes more visible and
the coefficients should be regarded as only approximately decorrelated
rather than strictly independent.
Figure~\ref{fig:wdm_orthogonality_decorrelation} illustrates these points
numerically for a small sampled transform and for white and coloured
noise realizations. The practical payoff of the (approximate) orthonormality and
decorrelation discussed above is that standard gravitational-wave
data analysis quantities --- the likelihood, signal-to-noise ratio,
and Fisher information matrix --- can all be written as simple sums
over WDM pixels, in direct analogy with the familiar frequency-domain
expressions.




\bibliographystyle{plain}
\bibliography{bib}

\end{document}
